\chapter{Gauge Geometry and Photon Identification}
\label{ch:ch08_photon}

The question of this chapter is which charge is the photon, and what a physical photon still needs. It treats the projector as a
field and the connection it induces, the compact families at the third rung, alignment by a gauged subalgebra, the identification of the
electric charge $B_{\rm em}$, and the sources and masslessness of the photon. The charge identification is established; the normal-plane
charge $Q_A$ is anomalous with the matter of $P_*$ and is not the photon. The Maxwell term and the selecting potential are adopted, not
derived, and electroweak breaking is missing; the adopted potential opposes it.

% --- source: tex/s11_fields#0 ((head))


% --- source: tex/s11_fields#1 (The projector as a field)
\section{The projector as a field}
\label{sec:fields}

Up to this point a projector has been a choice made once and for all. A choice of this kind becomes a field as soon as it is
allowed to vary from point to point, and the selection problem then changes character: instead of asking which projector is
right we ask which configurations an action prefers. Before any action can be written, three things must be known: the space in
which the projector moves, the connection it induces on its own image, and the sign of the natural metric on its variations. All
three can be computed exactly, and each of them carries a surprise.

% --- source: tex/s11_fields#2 (Configuration space)
\subsection{Configuration space}

Let $\mathcal P_r$ be the set of orthogonal projectors $P$ on $S_3=\R^{512}$ of rank $r$ that commute with the interior Lorentz
action $\Sigma$. Differentiating $P^2=P$ gives $dP\,P+P\,dP=dP$, hence
\[
P\,dP\,P=0,\qquad (\id-P)\,dP\,(\id-P)=0 :
\]
a variation of a projector is purely off-diagonal with respect to the splitting $\operatorname{im}P\oplus\ker P$. The tangent space
is therefore
\[
T_P\mathcal P_r=\bigl\{S:\ S^{\!\top}=S,\ PSP=0,\ (\id-P)S(\id-P)=0,\ [S,\Sigma]=0\bigr\},
\]
the tangent space of the real Grassmannian cut down by the Lorentz condition.

\begin{proposition}[Configuration space]
\label{prop:config}
The Lorentz commutant $\mathcal C\subset\End_\R(S_3)$ is spanned by the $32768$ blades that commute with the six interior
bivectors. Its centre is $\{\id,\Omega\}\cong\C$ with $\Omega=\gamma^0\gamma^1\gamma^2\gamma^3$, $\Omega^2=-\id$, and since a
semisimple real algebra of dimension $2\cdot128^2$ with centre $\C$ is $M_{128}(\C)$, we have $\mathcal C\cong M_{128}(\C)$.
Consequently
\[
\mathcal P_r\cong\mathrm{Gr}_\C\bigl(\tfrac r4,128\bigr),\qquad
\dim_\R\mathcal P_r=2\,\frac r4\Bigl(128-\frac r4\Bigr),
\]
which is $3584$ at $r=64$ and $6144$ at $r=128$, against $\dim\mathrm{Gr}_\R(64,512)=28672$ for the unconstrained real Grassmannian.
\end{proposition}

\begin{proof}
A blade commutes with $e_ae_b$ exactly when it contains both or neither of $e_a,e_b$, so the blades commuting with the six interior bivectors are those
containing all or none of the four interior generators: $2\cdot2^{14}=32768$. They are generated by the fourteen external generators and $\Omega$, and the
central ones are $\id$ and $\Omega$; a central element has only central blades in its support, since conjugation by a blade multiplies each blade by $\pm1$.
$\Omega^2=-\id$ because $\Omega$ has one positive and three negative generators. Every $e_A$ satisfies $e_A^{\!\top}=\pm e_A$, so $\mathcal C$ is closed under
transposition, and $\Tr(x^{\!\top}x)>0$ for $x\neq0$ excludes nilpotent ideals: $\mathcal C$ is semisimple, its centre is a field, so it is simple, and a
simple $\C$-algebra of complex dimension $128^2$ is $M_{128}(\C)$. As a module, $\R^{512}$ is two copies of $\C^{128}$, so a projector of real rank $r$ in
$\mathcal C$ has complex rank $r/4$. Transposition is a positive involution acting on the centre as conjugation ($\Omega^{\!\top}=-\Omega$), hence conjugate
to the Hermitian adjoint, so orthogonal projectors correspond to Hermitian ones and $\mathcal P_r\cong\mathrm{Gr}_\C(\tfrac r4,128)$.
\end{proof}

\noindent Executable check: a7-1 (exact symbolic or integer arithmetic; \cref{app:methods}).

The complex structure that appears here is the interior volume $\Omega$, not the internal one $\omega_I$ of \cref{thm:i}. That is
not a contradiction. The commutant of the Lorentz action alone contains $e^{(1)}_2$, which commutes with the six interior bivectors
but anticommutes with the fourteen-index volume $\omega_I$, so $\omega_I$ is not central in $\mathcal C$; it becomes central only in
the smaller algebra that also commutes with the internal bivectors, which is the algebra of \cref{thm:i}. Each structure is central
in its own ambient algebra, which is the pattern described after \cref{prop:6plus8plus1}.

\begin{remark}[No blade is tangent]
\label{rem:noblade}
A count of tangent directions by blades fails, for a reason worth recording. Let $P=\prod_i\tfrac12(\id+X_i)$ be one of the
projectors of \cref{sec:projectors}, built from $k$ commuting involutions, and let $B$ be a blade anticommuting with at least one
$X_i$. Then $PBP=0$, but $(\id-P)B(\id-P)=(\id-P-P_\varepsilon)B$ with $P_\varepsilon=\prod_i\tfrac12(\id+\varepsilon_iX_i)$, which
vanishes only when $P+P_\varepsilon=\id$, that is only for $k=1$. There is also a rank argument: a blade is a signed permutation matrix,
invertible of rank $512$, while a tangent vector has the block form $\bigl(\begin{smallmatrix}0&A^{\!\top}\\A&0\end{smallmatrix}\bigr)$ and
rank at most $2r$. For $k\ge2$ no single blade is tangent, the tangent space is not spanned by blades, and a blade count
overestimates it, by a factor four at $r=64$ and two at $r=128$. The dimensions of \cref{prop:config} are obtained by imposing
both conditions sector by sector.
\end{remark}

% --- source: tex/s11_fields#3 (The sign of a variation)
\subsection{The sign of a variation}


The trace form on an orbit splits by the Cartan involution (\cref{thm:ghost}): compact directions carry positive and boost directions negative kinetic weight.

The group that the geometry supplies is generated by the internal bivectors, and it moves a projector along an orbit inside
$\mathcal P_r$. Whether a bivector moves $P$ to another orthogonal projector, and with what sign the natural trace form measures
the motion, is decided by the signature rule of \cref{thm:sig}.



\begin{example}[The paired-mode projector at one constituent]
\label{ex:ghostPi}
Take the rank-four projector $\Pi$ of \cref{ex:pairedmode}. Its stabiliser in $\so{3,3}$ has dimension nine, so its orbit has
dimension six. On the bivectors
\[
B_I=\tfrac12\bigl(e_1e_2,\ e_1e_3,\ e_2e_3,\ e_1e_5,\ e_1e_6,\ e_2e_6\bigr),\qquad T_I=[B_I,\Pi],
\]
an exact $8\times8$ computation gives $\Tr(T_IT_J)=\operatorname{diag}(1,1,1,-1,-1,-1)$. The first three $B_I$ are time--time rotations,
the last three are boosts, exactly as \cref{thm:ghost} predicts. Only three of the fifteen basis bivectors commute with $\Pi$
individually, while the stabiliser has dimension nine, so six of its directions are combinations of basis bivectors and no single
one of them. Reading a stabiliser off basis elements overestimates the orbit.
\end{example}

Positive norm and tangency are nevertheless different conditions. For unit $T_{\rm sym}$ and $T_{\rm anti}$ of the two kinds,
$\Tr\bigl((T_{\rm sym}+\tfrac12T_{\rm anti})^2\bigr)=1-\tfrac14>0$, yet the direction is not symmetric and leaves the orthogonal
projectors. What is exact is the orthogonal splitting, not an equivalence between the sign of the norm and tangency. On the manifold
of orthogonal projectors itself every variation is symmetric and the trace form is positive definite, so the indefiniteness of the
example is an artefact of moving $\Pi$ along the full $\mathrm{Spin}(3,3)$ orbit. This is a statement about the target of a
projector field and not about a field theory over a base: in a base of split signature the kinetic contraction
$g^{\mu\nu}\Tr(D_\mu P\,D_\nu P)$ is indefinite in the spacetime index whatever the target metric is.

At three constituents the same reasoning counts the physical motions of $\Pfam$. Of the $91$ internal bivectors $43$ are compact,
forming the maximal compact $\so8\oplus\so6$, and the compact part of the stabiliser of $\Pfam$ has dimension nine
(\cref{thm:table}), so the compact orbit through $\Pfam$ has dimension
\[
43-9=34 ,
\]
which is also the rank of $B\mapsto[B,\Pfam]$ on the compact bivectors. Only seven of the $43$ compact basis bivectors commute with
$\Pfam$ individually, which repeats the warning of \cref{ex:ghostPi}. Of the $3584$ configuration directions at rank $64$,
thirty-four lie along the only group the geometry supplies. Whether these thirty-four are gauge redundancies or physical modes is a
question about an action, which the algebra does not supply.

% --- source: tex/s11_fields#4 (The induced connection over-generates)
\subsection{The induced connection over-generates}

A varying projector carries its own connection. If the gauge field is not postulated but induced, $\nabla=P\,d$ on the bundle of
images, then its curvature $F_{\mu\nu}=P[\partial_\mu P,\partial_\nu P]P$ is a composite, and the algebra in which it takes values is
computed rather than chosen. To leading order in the Ambrose--Singer theorem that algebra is
\[
\mathfrak h(P)=\operatorname{span}\bigl\{P[S,S']P:\ S,S'\in T_P\mathcal P_r\bigr\}.
\]

\begin{theorem}[Over-generation]
\label{thm:holonomy}
Let $P(x)$ range over all of $\mathcal P_{64}$. The curvature values $P[\partial_\mu P,\partial_\nu P]P$ are antisymmetric, hence
compact, and lie in the block, and their span on $\operatorname{im}\Pfam$ has dimension
\[
\dim\mathfrak h(\Pfam)=256=\dim\uu{16},
\]
the largest compact algebra the block admits at all. It exceeds the six dimensions of $\keff(\Pfam)=\uu1^3\oplus\su2$ by a factor
of about $43$.
\end{theorem}

\begin{proof}
Write $\Pfam=\prod_i\tfrac12(\id+X_i)$ with the three involutions of \cref{sec:projectors}. A tangent vector at $\Pfam$ is $S=Y+Y^{\!\top}$ with
$Y\in\Pfam\mathcal C(\id-\Pfam)$, so $T_{\Pfam}\mathcal P_{64}$ is spanned by the $S_b=\Pfam b(\id-\Pfam)+(\id-\Pfam)b^{\!\top}\Pfam$ for blades $b\in\mathcal C$, and
$\mathfrak h(\Pfam)$ by the $\Pfam[S_{b_1},S_{b_2}]\Pfam=\Pfam b_1(\id-\Pfam)b_2^{\!\top}\Pfam-\Pfam b_2(\id-\Pfam)b_1^{\!\top}\Pfam$. These are antisymmetric, lie in
the block and in $\mathcal C$, so $\mathfrak h(\Pfam)$ lies in the antisymmetric part of $\Pfam\mathcal C\Pfam$. A blade of $\mathcal C$ that anticommutes
with some $X_i$ has $\Pfam b\Pfam=0$ (\cref{rem:noblade}); a blade commuting with all three has $\Pfam b\Pfam=b\Pfam\neq0$, and $b$ and $bX_i$ give the same
element. The $4096$ blades of $\mathcal C$ commuting with $X_1,X_2,X_3$ fall into $512$ classes modulo $\langle X_1,X_2,X_3\rangle$; distinct classes are
trace-orthogonal, since a non-scalar blade is traceless, so they form a basis of $\Pfam\mathcal C\Pfam\cong M_{16}(\C)$; the sign in $b^{\!\top}=\pm b$ is
constant on a class, and $256$ classes are antisymmetric. Hence $\dim\mathfrak h(\Pfam)\le256$. Conversely, let $c$ be an antisymmetric class, $b_1=E_2E_7$,
which lies in $\mathcal C$ and anticommutes with $X_1$, and $b_2=b_1c$ up to sign, which then anticommutes with $X_1$ as well. Then $\Pfam b_i\Pfam=0$, so
$\Pfam[S_{b_1},S_{b_2}]\Pfam=\Pfam(b_1b_2^{\!\top}-b_2b_1^{\!\top})\Pfam$; this is antisymmetric and proportional to $c\Pfam$, and the proportionality factor is
$\pm2$ because $c$ is antisymmetric. So every antisymmetric class is reached and $\dim\mathfrak h(\Pfam)=256$. The class count is a finite
computation on blade parities; check a7-2 performs it and also computes the $256$ curvature values from the definition, certifying their span by an exact
modular rank of their integer Gram matrix.
\end{proof}

\noindent Executable check: a7-2 (exact symbolic or integer arithmetic; \cref{app:methods}).

The induced connection is not flat, every curvature value is compact, and every value lies in the block. The route fails by
over-generating and by nothing else. In particular chirality is not among the casualties. The span in the theorem is taken over all projector
fields. One field configuration, for instance a solution of a variational model, has curvature values in a subspace that can be much
smaller. So the theorem shows that the geometry alone does not confine the induced curvature to $\keff(\Pfam)$; it does not show that the
curvature of every solution leaves $\keff(\Pfam)$.

\begin{proposition}[Chirality must be imposed, and costs nothing]
\label{prop:chiralpreserved}
Every projector of the family commutes with $\Gamma$. If a field $P(x)$ is required to satisfy $[P(x),\Gamma]=0$ everywhere, then
$[\partial_\mu P,\Gamma]=0$ and
\[
\bigl[P[\partial_\mu P,\partial_\nu P]P,\ \Gamma\bigr]=0 ,
\]
so the induced curvature preserves chirality and the holonomy lies in $\uu8\oplus\uu8$. That algebra is still far larger than
$\uu1^3\oplus\su2$.
\end{proposition}

\noindent The tangent vectors commuting with $\Gamma$ are those built from even blades, and the argument of the proof of \cref{thm:holonomy}, restricted to
the $128$ even antisymmetric classes, shows that the span is all of $\uu8\oplus\uu8$. Executable check: a7-2 (exact symbolic or integer arithmetic;
\cref{app:methods}).

The comparison of \cref{thm:holonomy} is between what the construction permits and what is wanted, not an impossibility theorem
about any particular field. A constant field has zero induced curvature, and a field compatible with $\Gamma$ has holonomy inside
$\uu8\oplus\uu8$. What the comparison does establish is the direction of the problem. A parent action must \emph{cut} the freedom of
the projector, by a constraint that shrinks $T_P\mathcal P$ before the curvature is formed or by a potential that freezes all but the
intended holonomy directions, rather than build a gauge field up from nothing. This reverses the usual motivation for composite
gauge fields. \Cref{fig:freedom} collects the numbers.

\begin{figure}[t]
\centering
\begin{tikzpicture}
\begin{axis}[xbar,xmode=log,log basis x=10,width=8.6cm,height=6.4cm,
  xmin=2,xmax=200000,
  ytick=data,
  symbolic y coords={keff,orbit,u8u8,u16,P64,GrR},
  yticklabels={{$\keff(\Pfam)=\uu1^3\oplus\su2$},{compact orbit of $\Pfam$},{curvature, chirality imposed},{curvature, unrestricted},{Lorentz-compatible projectors $\mathcal P_{64}$},{real Grassmannian $\mathrm{Gr}_\R(64,512)$}},
  yticklabel style={font=\small},
  xlabel={real dimension (logarithmic)},
  bar width=9pt,
  nodes near coords,
  nodes near coords style={font=\small,anchor=west},
  point meta=rawx,
  enlarge y limits=0.12,
  axis x line*=bottom, axis y line*=left]
\addplot[fill=tcol!45,draw=tcol] coordinates {(6,keff) (34,orbit) (128,u8u8) (256,u16) (3584,P64) (28672,GrR)};
\end{axis}
\end{tikzpicture}
\caption{How much freedom a projector field has at rank $64$, and how little of it is wanted. The configuration space of
Lorentz-compatible projectors (\cref{prop:config}) is far smaller than the real Grassmannian but still enormous; the group the
geometry supplies moves $\Pfam$ along only $34$ directions; and the curvature algebra the induced connection can reach
(\cref{thm:holonomy,prop:chiralpreserved}) exceeds the faithful gauge algebra of the block by more than an order of magnitude.}
\label{fig:freedom}
\end{figure}

% --- source: tex/s11_fields#5 (One explicit variational model)
\subsection{One explicit variational model}
\label{sec:trialaction}

A list of fields does not determine field equations. What can be done without a derived action is to adopt one explicitly, derive
its equation, and record what it can and cannot do. On a four-dimensional Lorentzian base, with $D_\mu P=\partial_\mu P+[A_\mu,P]$ and
with $f$ and $V$ as inputs, take
\begin{equation}
\label{eq:trialS}
S_P=\int\sqrt{-g}\,\Bigl(\tfrac{f^2}2\Tr\bigl(D_\mu P\,D^\mu P\bigr)-V(P)\Bigr).
\end{equation}
Varying along the Grassmannian, $\delta P=[\epsilon,P]$ with $\epsilon$ antisymmetric, gives
\begin{equation}
\label{eq:Pieom}
\bigl[P,\ f^2D_\mu D^\mu P+V'(P)\bigr]=0 ,
\end{equation}
and for $V=0$ this is the conservation law $D_\mu\mathcal C^\mu=0$ of the matrix current $\mathcal C_\mu=[P,D_\mu P]$. The variation
has to be projected onto the admissible tangent space whenever $P$ is restricted further, for instance to commute with $\Gamma$.

Three consequences are immediate. First, a potential that depends only on the spectrum of a projector of fixed rank is constant,
since $\Tr f(P)=rf(1)+(d-r)f(0)$; such a potential selects nothing, and for the pure kinetic term every constant admissible projector
is a critical point. Selection needs couplings that distinguish orbits, and there is more than one orbit to distinguish: at rank
$64$ the orbits through $\Pfam$ and $P_*$ have dimensions $34$ and $20$ (\cref{prop:orbits}). Second, whether varying $A_\mu$ gives a
Yang--Mills equation depends on how the action depends on $A$, not on $A$ appearing in a list of fields; a composite $A$ contributes
through the chain rule to the equations of its constituents. Third, the pointwise stabiliser condition $[A_\mu,P]=0$ is not the
condition $D_\mu P=0$ that a moving subbundle be preserved; the two agree only for constant $P$. Which connection acts on matter,
$P(d+A)$, an ambient connection preserving $P$, or an independent connection in a moving frame of the block, is a further choice,
and we make none of them here.

% --- source: tex/s11_fields#6 (What Yang--Mills theory keeps open)
\subsection{What Yang--Mills theory keeps open}
\label{sec:ym}

Four conditions on a connection have been established as theorems: commutation with the Lorentz action
(\cref{thm:centraliser}); commutation with the projector, with a faithful action on its image (\cref{def:keff,thm:table});
antisymmetry, which for bivectors is \cref{thm:sig}; and preservation of the kinetic kernels $K^\mu$. They are homogeneous and
geometric, and together they determine the algebra in which $A_\mu$ takes its values, $\uu1^3\oplus\su2$ on the family block,
$\so6\oplus\su2^{\oplus3}$ at rank $256$ and the Pati--Salam algebra on $P_*$, over the single frame of \cref{prop:frame}. They do not
determine the field, since no parent action is derived; nor the couplings, since homogeneous conditions cancel every
normalisation; nor which abelian direction is electromagnetic, by \cref{thm:pairing}.

There is no representational obstruction behind this. By the theorem on universal connections \cite{NarasimhanRamanan}, every
connection on a bundle over a manifold of bounded dimension is the pullback of the canonical connection of a tautological bundle
over a Grassmannian of sufficient rank, so the projector representation of a gauge field is always available in principle. The
difficulty is dynamical, and power counting makes it concrete. A projector is dimensionless, so $\Tr(\partial P\,\partial P)$
carries a coupling of mass dimension $D-2$ and $\Tr F^2$ one of dimension $D-4$. In two dimensions the sigma model is
renormalisable and the composite gauge field acquires a kinetic term at leading order in $1/N$, with linear confinement of the
constituents \cite{DAddaDiVecchiaLuscher,Witten79}. In four dimensions the sigma model is not power-counting renormalisable and that
argument does not transfer.

Confinement is a separate question again. Invariant tensors permit meson-like and baryon-like operators, but they say nothing
about masses or about the energy of separated sources, and the diagnostics with dynamical matter are flux-tube and string-breaking
correlation matrices \cite{Bali05}. Moreover a block supplies a multiplicity index, not three constituents: three copies of a doublet
form a direct sum, whereas three constituents require a tensor product with antisymmetrisation, so even the operator statement needs
a Fock construction first.

% --- source: photon_selection/p1_families#0 ((head))
% New references: sections/bib_p1.tex.

% --- source: photon_selection/p1_families#1 (The compact families at the third rung)
\section{The compact families at the third rung}
\label{p1:sec:main}

The second-rung vacuum of \cref{def:towervac} lies on a compact family $S^4$ of Lorentz-invariant projectors, and a spectrum-only
potential is constant on it \cref{sec:photon}. A coupling that distinguishes its points has to be found at $n\ge3$, where colour survives
\cref{sec:assessment}. This section fixes the geometry on which any such coupling acts. The results are as follows. At the third rung the
Lorentz-invariant symmetric involutions that are products of four internal generators form three orbits of the compact group $K$. Only the
four-positive orbit $F_3$, an oriented Grassmannian of dimension $16$, commutes with colour, and it commutes with colour at every point
(fixed). The second-rung family is one Clifford four-sphere inside $F_3$ and is not $K$-invariant (fixed). The vacuum $\Pi_A$ built on the
auxiliary four-plane is the analogue of \cref{def:towervac}. It is neutral under the construction's charge and under electromagnetism, and
it is fixed by the Pati--Salam and Standard-Model algebras (fixed). Its sixteen fluctuations are four complex Standard-Model doublets of
hypercharge $\tfrac12$ (fixed). Any $K$-invariant coupling is constant on $F_3$. The alignment energy of a proper gauged subalgebra vanishes
exactly on its fixed set, and these fixed sets are determined here. A $K$-equivariant duality $P\mapsto P^\perp$ preserves every gauge-boson
mass matrix, so no alignment energy separates $\Pi_A$ from $\Pi_W$ (fixed). On the matter block the construction's charge acts as half the
fermion number (fixed), and with $\su2_L$ it is anomalous there (restated). Identifiers p1-$n$ refer to the \texttt{PASS} lines of
\texttt{p1\_families.py} (\cref{p1:tab:claims}).

\paragraph{Conventions.} The realisation is that of \cref{app:realisation}: $E_{6(k-1)+a}=\omega^{\otimes(k-1)}\otimes e_a\otimes\id^{\otimes(n-k)}$ on
$\R^{8^n}$, every generator and blade a signed permutation matrix, and $e_a^{\!\top}=\eta_{aa}e_a$. Products, commutators and traces of blades are therefore integer
computations, and blades obey the commutation rule, equation~\cref{eq:bladecomm}, and the transpose rule, \eqref{eq:transpose}.

\paragraph{Definitions.} Rung $n=3$, spectator interior action: interior generators $\gamma^\mu=(E_1,E_4,E_5,E_6)$; internal generators the other fourteen. The
positive internal span is $V_+=\operatorname{span}(A\cup W)$, of dimension $8$, and the negative internal span is $V_-=\operatorname{span}(C)$, of dimension $6$. The groups are those of \cref{sec:pstar}: auxiliary $A=(E_2,E_3,E_8,E_9)$,
weak $W=(E_7,E_{13},E_{14},E_{15})$, written $(w_1,\dots,w_4)$, and colour $C=(E_{10},E_{11},E_{12},E_{16},E_{17},E_{18})$.
$K=\mathrm{Spin}(8)\times\mathrm{Spin}(6)$ is the group generated by the $43$ compact internal bivectors \cref{thm:sig}, with Lie algebra $\mathfrak k=\so8\oplus\so6$. For an
oriented four-plane $P$ in $V_+$ with oriented orthonormal basis $(u_1,u_2,u_3,u_4)$, $\mathrm{vol}_P:=u_1u_2u_3u_4$, a symmetric involution commuting with every interior
gamma, and $\Pi_P:=\tfrac12(\id+\mathrm{vol}_P)$. The compact family $F_3:=\{\Pi_P:P\text{ an oriented four-plane in }V_+\}$ is the $K$-orbit of $\Pi_A$ ($A$ the
auxiliary four-plane) and is diffeomorphic to the oriented Grassmannian of four-planes in $\R^8$, of dimension $16$ (\cref{p1:thm:classify}). The gauged subalgebras considered
are: $\mathfrak h_{43}=\mathfrak k$, all compact internal bivectors, as in \cref{prop:alignment}; $\mathfrak h_{\rm PS}=\so4_W\oplus\so6_C$, the Pati--Salam algebra
$\su2_L\oplus\su2_R\oplus\su4$; $\mathfrak h_{\rm SM}=\su3_C\oplus\su2_L\oplus\uu1_Y$, embedded as in \cref{sec:pstar}; $\uu1_{\rm em}$, the bivector part of
$Q=-JT^3_L+Y$; and $\uu1_A=\operatorname{span}(q_A)$, the construction's charge $Q=\tfrac12J_1$ at $n=3$. The alignment energy of a gauged subalgebra $\mathfrak h$ with
orthonormal basis $(t_a)$ and couplings $g_a$ is $V_{\mathfrak h}(\Pi):=\sum_ag_a^2\|[t_a,\Pi]\|^2$ (Frobenius norm on $\R^{512}$). With the projector kinetic term
$\tfrac{f^2}2\Tr(D\Pi D\Pi)$ of \cref{sec:photon} it is proportional to the trace of the gauge-boson mass matrix at $\Pi$.

\paragraph{Notation.} $J$ is the ordered product of the internal generators, $\Gamma=J\gamma^0\gamma^1\gamma^2\gamma^3$, $q_A=E_2E_3$, $Q_A=\tfrac12q_A$, $\chi_W=w_1w_2w_3w_4$,
$\chi_C$ the product of $C$, and $P_*=\tfrac18(\id+\mathcal X)(\id+\mathcal Z)(\id+\mathcal U)$ with $\mathcal X=\Gamma\chi_W$, $\mathcal Z=-J\chi_C$, $\mathcal U=-Jq_A$ \cref{thm:PS}; $P_{L,R}=P_*\tfrac12(\id\mp\Gamma)$.
The colour Cartan element is $H_C=-(E_{10}E_{16}+E_{11}E_{17}+E_{12}E_{18})$, and $\su3_C$ is its derived centraliser in $\so6_C$ \cref{prop:PSledger}. Further
$\tau_i=-\tfrac12w_jw_k$ for cyclic $(i,j,k)$, $t^{L,R}_i=\tfrac12(\id\mp\chi_W)\tau_i$ \cref{thm:PSalg}, the hypercharge bivector is $t_Y=t^R_3+\tfrac16H_C$, and
\[
B_{\rm em}=\tau_3+\tfrac16H_C=t^L_3+t_Y=-\tfrac12E_7E_{13}+\tfrac16H_C
\]
spans $\uu1_{\rm em}$: on the block $Q=-J(t^3_L+t^3_R+\tfrac16H_C)$ \cref{prop:PSledger}. The vector representation $\rho$ is defined by $[X,u]=\rho(X)u$ for internal
vectors $u$. Thus $\rho(\tfrac12e_ie_j)=E_{ij}-E_{ji}$ on $V_+$, with a sign $-1$ on $V_-$, and the lift $M\mapsto\sum_{i<j}M_{ij}\tfrac12e_ie_j$ (sign $-1$ on $V_-$) is a Lie
algebra map inverse to $\rho$ (p1-12). We write $\rho_+(X)\in\so8$ for the $V_+$ part, $p_P$ for the orthogonal projector of $\R^8$ onto $P$, $\theta$ for $+1$ on $V_+$ and
$-1$ on $V_-$, and $\widetilde{\mathrm{Gr}}_k(V)$ for the oriented Grassmannian. On vectors, $\rho(Q_A)$ maps $E_2\mapsto-E_3$, $E_3\mapsto E_2$ and kills $E_8,E_9,W$. $\rho(B_{\rm em})$
maps $w_1\mapsto w_2$, $w_2\mapsto-w_1$ and kills $A,w_3,w_4$. $J_Y=2\rho(t_Y)$ is a complex structure on $W$ that kills $A$.

% --- source: photon_selection/p1_families#2 (The second rung, restated)
\subsection{The second rung, restated}
\label{p1:sec:n2}

\begin{proposition}[The second-rung family is a $K_2$-orbit; restated and fixed]
\label{p1:prop:n2}
At $n=2$ the internal complement of the spectator action has signature $(5,3)$, with $V_+^{(2)}=\operatorname{span}(E_2,E_3,E_7,E_8,E_9)$ and
$V_-^{(2)}=\operatorname{span}(E_{10},E_{11},E_{12})$. Let $(a,b,c,d,e)=(E_2,E_3,E_7,E_8,E_9)$, $Q=\tfrac12ab$, $H=abcd$, $\Pi_0=\tfrac12(\id+H)$, $I_5=abcde$, $Z_i=I_5v_i$.
\begin{enumerate}[label=(\roman*),leftmargin=2.2em]
\item \emph{(Restated.)} \crefrange{thm:towervac}{thm:towercurved} hold: $H$ is a symmetric traceless involution commuting with the interior gammas, $Q$ and the
internal volume; $[\tfrac12be,\Pi_0]=-\tfrac12acde$; the $Z_i$ are a symmetric Clifford system, $T_i=[\tfrac12v_ie,\Pi_0]=\tfrac12Z_i$ with $\Tr(T_iT_j)=16\delta_{ij}$, $Q$
rotates $(t_a,t_b)$ and fixes $t_c,t_d$; along $\Pi(n)=\tfrac12(\id+I_5n)$, $\Tr(d\Pi^2)=16|dn|^2$, $-\Tr([Q,\Pi]d\Pi)=16(n_adn_b-n_bdn_a)$ and $[\Pi,d\Pi]=\tfrac14[n,dn]$.
The stabiliser of $H$ in $\so{5,3}$ is $\so P\oplus\so{P^\perp}=\so4\oplus\so{1,3}$, of dimension $12$ with compact part $9$, and every other internal bivector
anticommutes with $H$ \cref{prop:towermatter}.
\item \emph{(Fixed.)} The $13$ compact internal bivectors preserve $\operatorname{span}\{Z_i\}$: $\so5$ acts on it by the vector representation and $\so3$
trivially. The $K_2$-orbit of $\Pi_0$, $K_2=\mathrm{Spin}(5)\times\mathrm{Spin}(3)$, is the whole family $S^4=\{\Pi(n)\}$, and
$\exp(\tfrac\vartheta2v_ie)$ carries $\Pi(e)$ to $\Pi(\cos\vartheta\,e+\sin\vartheta\,v_i)$. Moreover $I_5n=\mathrm{vol}_{n^\perp}$, so $S^4=\{\Pi_P:P\in\widetilde{\mathrm{Gr}}_4(V_+^{(2)})\}$.
\item \emph{(Fixed.)} $[Q,\Pi(n)]=0$ exactly on the two-sphere $n_a=n_b=0$, the planes containing $\operatorname{span}(a,b)$. Every $\Pi(n)$ has spectrum
$\{0^{32},1^{32}\}$, and $\tfrac1{64}\sum\|[e_ie_j,\Pi(n)]\|^2=4$ over the $13$ compact blades. A spectrum-only potential and the full alignment energy are
constant on $S^4$.
\end{enumerate}
\end{proposition}

\begin{proof}
(i) is the computation of \cref{sec:photon}, repeated in p1-1 to p1-4 and p1-7. The stabiliser statement is equation~\cref{eq:bladecomm}: a bivector $e_ie_j$ and the grade-four blade $H$ commute iff $|\{i,j\}\cap\{a,b,c,d\}|$ is even. (ii) A positive bivector inside $V_+^{(2)}$ shares two
indices with $I_5$ and commutes with it, so $[X,Z_i]=I_5[X,v_i]=I_5\rho(X)v_i$; a negative bivector shares none and commutes with each $Z_i$. Since $\mathrm{Spin}(5)$ is
transitive on $S^4$, the orbit of $\Pi(e)$ is the family. The rotation formula is $\Ad(\exp\tfrac\vartheta2v_ie)e=\cos\vartheta\,e+\sin\vartheta\,v_i$.
If $(f_1,\dots,f_4,n)$ is an oriented orthonormal basis of $V_+^{(2)}$, then $I_5=f_1f_2f_3f_4n$ and $I_5n=f_1f_2f_3f_4$ (p1-5). (iii) $[Q,Z_i]=I_5\rho(Q)v_i$, and
$\rho(Q)$ rotates $(a,b)$ and kills $c,d,e$ (p1-6). Each $\Pi(n)$ is conjugate to $\Pi_0$ by an orthogonal matrix, and the alignment sum is $K_2$-invariant, as in
\cref{p1:prop:K} below (p1-4, p1-5).
\end{proof}

At $n=2$ three descriptions coincide: the Clifford sphere of five positive generators, the $K_2$-orbit, and the oriented Grassmannian of four-planes in
$V_+^{(2)}$. At $n=3$ they separate. The Grassmannian description is the one that continues.

% --- source: photon_selection/p1_families#3 (Four-generator involutions at the third rung)
\subsection{Four-generator involutions at the third rung}
\label{p1:sec:classify}

\begin{lemma}[Blade census; fixed]
\label{p1:lem:census}
At $n=3$ the eighteen generators satisfy the Clifford relations of $\R^{9,9}$, $V_+=\operatorname{span}(E_2,E_3,E_7,E_8,E_9,E_{13},E_{14},E_{15})$, and the $43$ compact
internal bivectors square to $-\id$ while the $48$ mixed ones square to $+\id$ (p1-8).
\begin{enumerate}[label=(\roman*),leftmargin=2.2em]
\item Of the $3060$ blades of grade four exactly the $1001$ internal ones commute with the four interior gammas. All are traceless. One with $q$ negative factors is
a symmetric involution if $q$ is even and an antisymmetric complex structure if $q$ is odd. The counts are $(4,0):70$, $(3,1):336$, $(2,2):420$, $(1,3):160$, $(0,4):15$
(p1-9).
\item More generally an internal blade of even grade $r$ with $q$ negative factors is a symmetric involution iff $r(r-1)/2+q$ is even. The only nontrivial one invariant
under $K$ is $\mathrm{vol}_+=E_2E_3E_7E_8E_9E_{13}E_{14}E_{15}$, and $\mathrm{vol}_+=-\mathcal Z=J\chi_C$ (p1-10).
\end{enumerate}
\end{lemma}

\begin{proof}
For $|I|=4$, equation~\cref{eq:bladecomm} gives $e_Ie_j=(-1)^{4-|I\cap\{j\}|}e_je_I$, so a generator commutes with $e_I$ iff it is not a factor. By equation~\cref{eq:transpose}, $e_I^2=(-1)^{r(r-1)/2+q}\id$ and $e_I^{\!\top}=e_I^2e_I$. An even internal blade commutes with $e_ie_j$ iff $|I\cap\{i,j\}|$ is even. For all same-sign pairs this
forces $I\cap V_\pm\in\{\emptyset,V_\pm\}$, which leaves $\mathrm{vol}_+$, $\chi_C$ ($r=q=6$, square $-\id$) and $\pm J$ ($r=14$, square $-\id$). The
identity $\mathrm{vol}_+=J\chi_C$ is checked on the matrices.
\end{proof}

\begin{lemma}[Transpose and recovery; fixed]
\label{p1:lem:transpose}
Let $P$ be a nondegenerate internal four-plane with $\eta$-orthonormal oriented basis $(u_1,\dots,u_4)$ and $\mathrm{vol}_P=u_1u_2u_3u_4$. Then $\mathrm{vol}_P$ commutes with the
interior gammas, $\mathrm{vol}_P^2=(-1)^q\id$, the internal vectors anticommuting with $\mathrm{vol}_P$ are exactly those of $P$, and
\[
\mathrm{vol}_P^{\!\top}=\mathrm{vol}_{\theta P}.
\]
Hence for $q$ even, $\Pi_P$ is an orthogonal projector iff $\theta P=P$, that is iff $P=(P\cap V_+)\oplus(P\cap V_-)$. For a positive-definite $P\not\subset V_+$, $\Pi_P$ is
an idempotent but not symmetric (p1-13).
\end{lemma}

\begin{proof}
The interior gammas anticommute with every internal vector, hence commute with a product of four. By equation~\cref{eq:transpose}, $\mathrm{vol}_P^2=(-1)^6\prod\eta(u_i)\id=(-1)^q\id$.
For an internal vector $u$, $u^{\!\top}=\theta u$. Transposing reverses the product, and reversing four anticommuting vectors costs $(-1)^{6}=1$. So
$\mathrm{vol}_P^{\!\top}=\theta u_1\cdots\theta u_4=\mathrm{vol}_{\theta P}$, and $\theta$ is an isometry of $\eta$. Each $u_i$ commutes with itself and anticommutes with the other three
factors, so every vector of $P$ anticommutes with $\mathrm{vol}_P$. A vector $\eta$-orthogonal to $P$ anticommutes with all four factors, so it commutes with $\mathrm{vol}_P$. Since $P$ is
nondegenerate, an internal $v$ splits as $v_P+v_\perp$; if $v$ anticommutes with $\mathrm{vol}_P$, so does $v_\perp$, which also commutes, hence $v_\perp\mathrm{vol}_P=0$ and $v_\perp=0$.
So $P$ is determined by $\mathrm{vol}_P$, and its orientation by the sign. Then $\mathrm{vol}_P$ is symmetric iff $\theta P=P$ as oriented planes.
\end{proof}

\begin{lemma}[Colour acts irreducibly on $V_-$; fixed]
\label{p1:lem:colour}
The centraliser of $\rho(H_C)$ in $\so{V_-}$ is $\uu3$. Its part orthogonal to $H_C$ is $\su3_C$, of dimension $8$ and perfect. Lifted to $\mathfrak k$, $\su3_C$ commutes
with $H_C$ and acts trivially on $V_+$. Its commutant on $V_-=\R^6$ is $\operatorname{span}(\id,\tfrac12\rho(H_C))$ with $(\tfrac12\rho(H_C))^2=-\id$, a field, so
$\su3_C$ acts irreducibly on $V_-$ (p1-15).
\end{lemma}

\begin{proof}
The computation is in the six-dimensional vector model and is lifted by the Lie algebra map of p1-12. A proper invariant subspace of an orthogonal representation
has an invariant orthogonal complement, and the projector onto it is an idempotent in the commutant other than $0$ and $\id$. A field has no such idempotent.
\end{proof}

\begin{theorem}[The three orbit types and the colour test; fixed]
\label{p1:thm:classify}
The Lorentz-invariant symmetric involutions $\mathrm{vol}_P$, with $P$ a nondegenerate internal four-plane, form three $K$-orbits. Their stabilisers and dimensions are:
\begin{center}
\footnotesize
\begin{tabular}{@{}lllccc@{}}
\toprule
type & orbit & stabiliser in $\so{8,6}$ & dim & compact part & orbit dim\\
\midrule
$(4,0)$ & $F_3\cong\widetilde{\mathrm{Gr}}_4(V_+)$ & $\so4\oplus\so{4,6}$ & $51$ & $27$: $\so4_A\oplus\so4_W\oplus\so6_C$ & $16$\\
$(2,2)$ & $\widetilde{\mathrm{Gr}}_2(V_+)\times\widetilde{\mathrm{Gr}}_2(V_-)/\pm$ & $\so{2,2}\oplus\so{6,4}$ & $51$ & $23$ & $20$\\
$(0,4)$ & $\widetilde{\mathrm{Gr}}_4(V_-)$ & $\so4\oplus\so{8,2}$ & $51$ & $35$ & $8$\\
\bottomrule
\end{tabular}
\end{center}
Every point has spectrum $\{+1^{256},-1^{256}\}$, so no spectrum-only potential separates the points or the types. Every point of $F_3$ commutes with all of
$\so6_C\supset\su3_C$. No point of the $(2,2)$ or $(0,4)$ orbits commutes with $\su3_C$.
\end{theorem}

\begin{proof}
By \cref{p1:lem:transpose} the symmetric ones are the $\mathrm{vol}_{P_+}\mathrm{vol}_{P_-}$ with $P_\pm\subset V_\pm$ of dimensions $(p,q)$, $q$ even. For $g\in K$,
$\Ad(g)u=\rho(g)u$ on vectors, so $\Ad(g)\mathrm{vol}_P=\mathrm{vol}_{\rho(g)P}$ (p1-14). $SO(8)\times SO(6)$ is transitive on pairs of oriented planes, which gives one orbit per
type; reversing both orientations in type $(2,2)$ gives the same element. By the recovery statement of \cref{p1:lem:transpose}, $\exp(tX)$ fixes $\mathrm{vol}_P$ iff $\rho(X)$
preserves $P$, so the stabiliser is $\so P\oplus\so{P^\perp}$. Bivector by bivector, $e_ie_j$ commutes with $\mathrm{vol}_P$ iff both or neither of $i,j$ lie in $P$, and
anticommutes otherwise; the counts are p1-11. The orbit dimension is $43$ minus the compact part. A symmetric involution of trace zero has eigenvalues $\pm1$, each $256$
times (p1-42). A vector of $V_+$ anticommutes with both factors of a colour bivector and so commutes with it; hence every colour bivector commutes with every $\mathrm{vol}_P$, $P\subset V_+$. For types $(2,2)$ and $(0,4)$,
$\su3_C$ fixes $\mathrm{vol}_P$ iff it preserves the plane $P_-\subset V_-$ of dimension $2$ or $4$, which \cref{p1:lem:colour} excludes. p1-16 checks this at the
representatives and at random $K$-translates.
\end{proof}

At $n=3$ the stabiliser of the four-positive type is the row of \cref{prop:towermatter}. The same proof covers all even grades: a $\theta$-stable plane of type $(p,q)$
with $r(r-1)/2+q$ even gives a $K$-orbit of dimension $p(8-p)+q(6-q)$. Among them only $\mathrm{vol}_+=-\mathcal Z$ is $K$-fixed (\cref{p1:lem:census}). No compact generator moves
$\tfrac12(\id+\mathrm{vol}_+)$, as for the chirality projectors of \cref{prop:rigid}; only the boosts do. \Cref{p1:lem:transpose} agrees with \cref{thm:ghost}, by which a
variation $[B,\Pi]$ stays among the orthogonal projectors exactly when $B$ is compact: a noncompact generator tilts $P$ out of $V_+$, and $\Pi_P$ stops being symmetric.

\begin{proposition}[Clifford spheres in $F_3$; fixed]
\label{p1:prop:spheres}
\begin{enumerate}[label=(\roman*),leftmargin=2.2em]
\item For an oriented five-plane $L\subset V_+$ with oriented orthonormal basis $(v_1,\dots,v_5)$ put $I_L=v_1\cdots v_5$ and $Z_i=I_Lv_i$. The $Z_i$ are a symmetric Clifford
system, and $\tfrac12(\id+I_Ln)=\Pi_{n^\perp\cap L}$. Hence $S_L:=\{\tfrac12(\id+I_Ln):|n|=1\}=\{\Pi_P:P\subset L\}\cong S^4$ is the $\mathrm{Spin}(L)$-orbit of any of its points.
Dually, for an oriented three-plane $U$ with volume $U_3$, $S^U:=\{\Pi_P:P\supset U\}=\{\tfrac12(\id+nU_3):n\in U^\perp,|n|=1\}\cong S^4$.
\item Through each $\Pi_P$ pass the spheres $S_L$ with $L=P\oplus\R u$, $u\in P^\perp$ (a projective space $\R P^3$ of them), with tangent $P\otimes u$, and the spheres $S^U$
with $U=P\cap a^\perp$, $a\in P$, with tangent $a\otimes P^\perp$. These are the two kinds of maximal linear spaces of rank-one elements in $T_{\Pi_P}F_3\cong P\otimes P^\perp$. The
inner spheres are parametrised by $\mathrm{Gr}_5(\R^8)$, of dimension $15$, and none of them is $K$-invariant.
\item On the inner sphere through $A\oplus\R E_7$ with $Q=Q_A$, and on the outer sphere of planes containing $\operatorname{span}(E_3,E_8,E_9)$ with $Q=B_{\rm em}$,
$\Tr(d\Pi^2)=128|dn|^2$ and $-\Tr([Q,\Pi]d\Pi)=\pm128(n_adn_b-n_bdn_a)$. On the inner sphere $[\Pi,d\Pi]=\tfrac14[n,dn]$. This is the structure of
\cref{thm:S4,thm:towercurved} with $16$ replaced by $128$ (p1-19, p1-32).
\item The second-rung family embedded with its own generators, $L_0=\operatorname{span}(a,b,c,d,e)=\operatorname{span}(E_2,E_3,E_7,E_8,E_9)$, is the inner sphere $S_{L_0}$. Its
point $n=c$ is $\Pi_A$, since $I_5c=E_2E_3E_8E_9$, and its point $n=e$ is $\Pi_0=\tfrac12(\id+E_2E_3E_7E_8)$. On $S_{L_0}$, $Q_A=\tfrac12ab$ is unbroken on the two-sphere
$n_a=n_b=0$. The algebras $\uu1_{\rm em}$, $\mathfrak h_{\rm SM}$ and $\mathfrak h_{\rm PS}$ are each unbroken only at $n=\pm c$, that is at $\Pi_{\pm A}$ (p1-47).
\end{enumerate}
\end{proposition}

\begin{proof}
(i) As in \cref{p1:prop:n2}(ii): $Z_iZ_j+Z_jZ_i=2\delta_{ij}$ because $I_L$ commutes with each $v_i$ and squares to $+\id$, and $I_Ln=f_1f_2f_3f_4$ for an oriented basis
$(f_1,\dots,f_4,n)$ of $L$. For $S^U$, $nU_3$ is the volume of $\operatorname{span}(n,U)$; vectors of $U^\perp$ anticommute with $U_3$, so the $nU_3$ again form a Clifford system
(p1-19). (ii) The orbit of $\mathrm{Spin}(L)$, respectively $\mathrm{Spin}(U^\perp)$, has tangent $[\mathfrak m_P\cap\so L,\Pi_P]$, respectively $[\mathfrak m_P\cap\so{U^\perp},\Pi_P]$,
where $\mathfrak m_P=P\wedge P^\perp$ is spanned by the bivectors mixing $P$ and $P^\perp$. These intersections are $P\wedge u$ and $a\wedge P^\perp$. (iii) $\Pi$ and $d\Pi$ are
combinations of the $Z_i$, and $\Tr(Z_iZ_j)=512\delta_{ij}$; $[Q,Z_i]=\pm Z_j$ by the action of $\rho(Q)$ on $L$ or on $U^\perp$ (p1-32). (iv) The identities are blade
identities. The fixed points follow from \cref{p1:thm:fixed} below, since $P\subset L_0$ contains or is orthogonal to $\operatorname{span}(E_7,E_{13})$ only for $P=A$, and
$E_{13}\notin L_0$; p1-47 computes the kernels of $n\mapsto[X,I_5n]$ directly.
\end{proof}

\begin{proposition}[The duality $P\mapsto P^\perp$; fixed]
\label{p1:prop:duality}
$\mathrm{vol}_+$ commutes with $\mathfrak k$, with the interior gammas and with every $\mathrm{vol}_P$, and $\mathrm{vol}_+\mathrm{vol}_A=\mathrm{vol}_W$. The map
$D(\Pi)=\tfrac12\bigl(\id+\mathrm{vol}_+(2\Pi-\id)\bigr)$ sends $\Pi_P$ to $\Pi_{P^\perp}$, with $P^\perp\subset V_+$ suitably oriented. It is a $K$-equivariant involution of $F_3$
without fixed points, and $D(\Pi_A)=\Pi_W$. For every $T,T'\in\mathfrak k$,
\[
\Tr\bigl([T,D\Pi][T',D\Pi]\bigr)=\Tr\bigl([T,\Pi][T',\Pi]\bigr),\qquad \Tr\bigl([T,\id-\Pi][T',\id-\Pi]\bigr)=\Tr\bigl([T,\Pi][T',\Pi]\bigr).
\]
So every gauge-boson mass matrix, and every alignment energy $V_{\mathfrak h}$ with $\mathfrak h\subset\mathfrak k$, takes equal values at $\Pi_{\pm A}$ and $\Pi_{\pm W}$. On the
block, $\mathrm{vol}_+P_*=-P_*$, and $\Tr(\Pi_AP_L)=32$, $\Tr(\Pi_WP_L)=0$, $\Tr(\Pi_WP_R)=32$ (p1-20, p1-46).
\end{proposition}

\begin{proof}
$\mathrm{vol}_+$ is $K$-invariant by \cref{p1:lem:census}, and it commutes with the interior gammas, with which it shares no index. In an orthonormal basis $(u_1,\dots,u_8)$
of $V_+$ adapted to $P\oplus P^\perp$, $\mathrm{vol}_+=\pm u_1\cdots u_8$, so $\mathrm{vol}_+$ commutes with $\mathrm{vol}_P=u_1u_2u_3u_4$ and
$\mathrm{vol}_+\mathrm{vol}_P=\pm u_5u_6u_7u_8=\pm\mathrm{vol}_{P^\perp}$. $D$ is an involution because $\mathrm{vol}_+^2=\id$, and it has no fixed point because $P\neq P^\perp$. Since $\mathrm{vol}_+$
commutes with $T$ and $\Pi$, $[T,D\Pi]=\mathrm{vol}_+[T,\Pi]$, and $\mathrm{vol}_+^2=\id$ gives the identity. The second identity is $[T,\id-\Pi]=-[T,\Pi]$. On the block,
$\mathrm{vol}_+P_*=-ZP_*=-P_*$. The traces follow from \cref{p1:prop:block}.
\end{proof}

% --- source: photon_selection/p1_families#4 (The vacuum _APi\_A)
\subsection{The vacuum \texorpdfstring{$\Pi_A$}{Pi\_A}}
\label{p1:sec:vacuum}

\begin{definition}[The third-rung vacuum]
\label{p1:def:vacuum}
$\Pi_A:=\tfrac12(\id+E_2E_3E_8E_9)$, built from the surplus pairs of the first and second constituents, and $\Pi_W:=\tfrac12(\id+\chi_W)$. Their orientation
reversals are $\Pi_{-A}=\id-\Pi_A$ and $\Pi_{-W}=\id-\Pi_W$.
\end{definition}

\begin{theorem}[Lorentz invariant, neutral and not rigid; fixed]
\label{p1:thm:vacuum}
$\Pi_A$ and $\Pi_W$ are rank-$256$ orthogonal projectors with the following properties.
\begin{enumerate}[label=(\roman*),leftmargin=2.2em]
\item Each commutes with the interior gammas, with $J$, with $q_A$, with $B_{\rm em}$, with the $21$ generators of $\mathfrak h_{\rm PS}$ and with the $12$ of
$\mathfrak h_{\rm SM}$ (p1-21).
\item The compact internal bivectors that move $\Pi_A$ are exactly the $16$ that mix $A$ with $W$. The tangents $T_{aw}=[\tfrac12aw,\Pi_A]=\tfrac12aw\,\mathrm{vol}_A$ are
symmetric with $\Tr(T_{aw}T_{a'w'})=128\,\delta_{aa'}\delta_{ww'}$; the same holds at $\Pi_W$ (p1-23).
\item The stabiliser of $\Pi_A$ in $\mathfrak k$ is $\so4_A\oplus\so4_W\oplus\so6_C$, of dimension $27$. It contains $\uu1_A$, $\uu1_{\rm em}$, $\mathfrak h_{\rm PS}$ and
$\mathfrak h_{\rm SM}$ (p1-35).
\end{enumerate}
The literal embedding $\tfrac12(\id+E_2E_3E_7E_8)$ of the second-rung vacuum is also a point of $F_3$ and commutes with $q_A$. It does not commute with $B_{\rm em}$ or with
$\su2_L$, because $E_7\in W$ (p1-22).
\end{theorem}

\begin{proof}
All statements are equation~\cref{eq:bladecomm} applied to blades. $\mathrm{vol}_A$ shares two indices with $q_A$, none with $E_7E_{13}$, the colour bivectors or the
weak bivectors, and four with $J$. A bivector $aw$ shares one index with $\mathrm{vol}_A$ and anticommutes with it, so
$[\tfrac12aw,\Pi_A]=\tfrac14[aw,\mathrm{vol}_A]=\tfrac12aw\,\mathrm{vol}_A$. Symmetry and the trace follow from $(aw\,\mathrm{vol}_A)^2=-(aw)^2=\id$ and the tracelessness of
nontrivial blades. $B_{\rm em}$, $t^{L,R}_i$ and $\su3_C$ are combinations of bivectors inside $W$ or $C$.
\end{proof}

In \cref{prop:spectator} the compact spectator algebra carries colour on three negative planes, weak isospin on two positive planes and a family $\so4$ on the remaining
two. With the grouping of \cref{sec:pstar} the remaining two planes are those of $A$, so $\Pi_A$ separates the weak four-plane from the family four-plane and keeps
everything else (interpretive). By \cref{p1:prop:spheres}(iv), $\Pi_A$ is the point $n=c$ of the second-rung family itself, embedded at $n=3$. The point $n=e$ chosen in
\cref{def:towervac} is not neutral at $n=3$.

\begin{theorem}[The fluctuations at $\Pi_A$; fixed]
\label{p1:thm:tangent}
Let $T_{\Pi_A}F_3=\operatorname{span}\{T_{aw}\}\cong A\otimes W$.
\begin{enumerate}[label=(\roman*),leftmargin=2.2em]
\item Under $\so4_A\oplus\so4_W\cong\su2^{\oplus4}$ it is the real irreducible $(\mathbf2,\mathbf2;\mathbf2,\mathbf2)$, with commutant $\R$ and every $\su2$ Casimir equal to
$\tfrac34$ (p1-24).
\item Under $\mathfrak h_{\rm SM}$, $\su3_C$ acts trivially; $\su2_L$ has Casimir $\tfrac34$, and $\ad t_Y$ satisfies $(\ad t_Y)^2=-\tfrac14$ and commutes with $\su2_L$. The space is four
complex doublets $(\mathbf1,\mathbf2)_{1/2}$, one for each direction of $A$, and $\so4_A$ rotates the doublet index (p1-25).
\item Under $\uu1_{\rm em}$, $\ad B_{\rm em}=\ad(t^3_L+t_Y)$. The charged directions are $A\otimes\operatorname{span}(w_1,w_2)$, four complex scalars
$\phi_a=x_{a,13}+ix_{a,7}$ of charge one, the upper components of the doublets. The neutral directions are $A\otimes\operatorname{span}(w_3,w_4)$, eight real, the lower
components.
\item Under $\uu1_A$ the charged directions are $\operatorname{span}(E_2,E_3)\otimes W$, four complex scalars $\phi_w=x_{2,w}+ix_{3,w}$. The neutral directions are
$\operatorname{span}(E_8,E_9)\otimes W$, eight real.
\item The same holds at $\Pi_W$ with $\operatorname{Hom}(W,A)\cong W\otimes A$ (p1-24, p1-25, p1-26).
\end{enumerate}
Here $x_{a,w}$ are the coordinates of $\Pi=\Ad(e^X)\Pi_A$, $X=\sum x_{a,w}\tfrac12aw$. At $n=2$ the same count gives one charged complex scalar and two neutral real
scalars \cref{thm:S4}.
\end{theorem}

\begin{proof}
For $X$ in the stabiliser, $[X,[Y,\Pi_A]]=[[X,Y],\Pi_A]$. So $T_{\Pi_A}F_3$ is the representation of the stabiliser on $\mathfrak m_A=A\wedge W\cong A\otimes W$, on which
$\rho(\tfrac12aw)$ is the rotation of the $(a,w)$ plane. All five statements reduce to the action on $A$ and $W$ listed in the conventions. $\su2_L$ acts on $W$ with a
quaternionic commutant (p1-37), so $W$ is one complex doublet with complex structure $J_Y=2\rho(t_Y)$. $\rho(t^3_L+t_Y)=\rho(\tau_3)$ rotates $(w_1,w_2)$ at unit rate and
kills $(w_3,w_4)$. p1-26 gives the explicit pairs, for example $[Q_A,T_{2,w}]=-T_{3,w}$, $[Q_A,T_{3,w}]=T_{2,w}$ and $[B_{\rm em},T_{a,7}]=T_{a,13}$, $[B_{\rm em},T_{a,13}]=-T_{a,7}$.
\end{proof}

\begin{proposition}[The vacuum on the matter block; fixed, and restated]
\label{p1:prop:block}
\begin{enumerate}[label=(\roman*),leftmargin=2.2em]
\item $\mathrm{vol}_A$ and $\chi_W$ commute with $\mathcal X,\mathcal Z,\mathcal U$, and $\mathrm{vol}_AP_*=-\Gamma P_*$, $\chi_WP_*=\Gamma P_*$. Hence $\Pi_AP_*=P_L$ and $\Pi_WP_*=P_R$, the two Weyl blocks of rank
$32$ (p1-27).
\item $q_AP_*=JP_*$. So $Q_A=\tfrac12q_A$, the third-rung form of the construction's charge $\tfrac12J_1$, acts on the block as $\tfrac12J$: its charge operator
$-JQ_A$ equals $\tfrac12$ on all sixteen Weyl components. It is half the fermion number of \cref{sec:pstar}, not the electric charge, which takes the values
$-1,-\tfrac13,0,\tfrac23$ (p1-18, p1-27).
\item Every $T_{aw}$ anticommutes with $\mathcal X=\Gamma\chi_W$, and $P_*T_{aw}P_*=0$: the fluctuations of $\Pi_A$ map the block into its orthogonal complement (p1-27).
\item \emph{(Restated, \cref{prop:PSanom} with $a=\tfrac12$, $b=0$.)} With $\hat Q_A=-JQ_A$ and $\hat T=-Jt^L_3$, we have $\Tr(P_L\hat Q_A\hat T^2)=4$ and
$\Tr(P_R\hat Q_A\hat T^2)=0$, and likewise $4$ for $\su2_R$ on $P_R$. For $\hat Q_{\rm em}=-JB_{\rm em}$, $\Tr(P_L\hat Q_{\rm em}\hat T^2)=0$. The mixed $SU(2)_L^2\,U(1)_A$ anomaly does
not cancel on the block, while the gravitational one does (p1-28) \cite{BouchiatIliopoulosMeyer72,GeorgiGlashow72}.
\end{enumerate}
\end{proposition}

\begin{proof}
(i) $\mathrm{vol}_A=\mathrm{vol}_+\chi_W$ by \cref{p1:prop:duality}. $\mathrm{vol}_+$ commutes with $\Gamma$ and $\mathrm{vol}_+P_*=-P_*$; with $\chi_WP_*=\Gamma P_*$
\cref{thm:PS} this gives $\mathrm{vol}_AP_*=-\Gamma P_*$. (ii) $\mathcal UP_*=P_*$ means $-Jq_AP_*=P_*$, and $J^2=-\id$ gives $q_AP_*=JP_*$; then $-JQ_AP_*=\tfrac12P_*$. (iii)
$T_{aw}$ is half a blade of grade four with three factors in $A$ and one in $W$. It shares one index with $\chi_W$ and anticommutes with it, and it commutes with $\Gamma$,
which is $\pm$ the product of all eighteen generators. Hence it anticommutes with $\mathcal X$ (p1-27).
Since $\mathcal XP_*=P_*$, $P_*TP_*=P_*\tfrac12(\id+\mathcal X)TP_*=P_*T\tfrac12(\id-\mathcal X)P_*=0$. (iv) On $P_L$, $\hat Q_A=\tfrac12$ and $\hat T^2=-(t^L_3)^2=\tfrac14$, and $\Tr P_L=32$; $t^L_3P_R=0$. The electric charges of the left doublets sum to zero.
\end{proof}

The charge that the second-rung construction gauges, $Q=\tfrac12J_1$, therefore becomes at $n=3$ a universal charge on the
$P_*$ matter, and with $\su2_L$ it cannot be gauged on that matter. The charge for which $\Pi_A$ plays the role of \cref{def:towervac} with electromagnetic content is
$B_{\rm em}$. Both meet the massless criterion of \cref{cor:massless} at $\Pi_A$ (\cref{p1:thm:kinetic}). Which of them is the photon is an identification and is
not decided by the geometry of $F_3$ (open).

% --- source: photon_selection/p1_families#5 (Kinetic term, current and transport on F_3F\_3)
\subsection{Kinetic term, current and transport on \texorpdfstring{$F_3$}{F\_3}}
\label{p1:sec:kinetic}

\begin{lemma}[Vector model; fixed]
\label{p1:lem:vector}
For $P\in\widetilde{\mathrm{Gr}}_4(V_+)$ and $X,Y\in\mathfrak k$,
\[
\Tr\bigl([X,\Pi_P][Y,\Pi_P]\bigr)=64\,\tr_8\bigl([\rho_+X,p_P][\rho_+Y,p_P]\bigr).
\]
Hence on $F_3$: $\tfrac{f^2}2\Tr(D\Pi D\Pi)=32f^2\tr_8(Dp\,Dp)$; for $Q\in\mathfrak k$, $j_\Pi=-f^2\Tr([Q,\Pi]D\Pi)=-64f^2\tr_8([\rho_+Q,p]Dp)$; and
\[
V_{\mathfrak h}(\Pi_P)=64\sum_ag_a^2\,\|[\rho_+t_a,p_P]\|^2 ,
\]
in which the $\so6$ components of the $t_a$ drop out (p1-29, p1-30).
\end{lemma}

\begin{proof}
Both sides are bilinear in $(X,Y)$ and invariant under $K$, because $\Ad(g)\Pi_P=\Pi_{\rho(g)P}$, $p_{RP}=Rp_PR^{\!\top}$ and $\rho$ is equivariant. So it suffices to take
$P=A$. Write $X=X_0+\sum x_{aw}\tfrac12aw$ with $X_0$ in the stabiliser; then $[X,\Pi_A]=\sum x_{aw}T_{aw}$, and the left side is $128\sum x_{aw}y_{aw}$. On the right,
$[\rho(\tfrac12aw),p_A]=-(E_{aw}+E_{wa})$, and $\rho_+X_0$ commutes with $p_A$, so the right side is $64\cdot2\sum x_{aw}y_{aw}$. Every tangent vector at $\Pi_P$ is $[Y,\Pi_P]$ for
some $Y\in\mathfrak k$, which gives the kinetic and current forms.
\end{proof}

The factor $64=512/8$ is the real trace on $\R^{512}$ against that on $\R^8$. At $n=2$ the same factor is $8=64/8$. With $\Tr(T_iT_j)=16\delta_{ij}$ there, this reproduces
the $8f^2$ and $16f^2$ of \cref{thm:S4}.

\begin{theorem}[Kinetic term and current at $\Pi_A$; fixed]
\label{p1:thm:kinetic}
With $\Pi=\Ad(e^X)\Pi_A$, $X=\sum x_{aw}\tfrac12aw$, and the conventions of \cref{thm:S4}, to quadratic order in $x$:
\[
\mathcal L_\Pi=64f^2\sum_{a,w}D_\mu x_{aw}D^\mu x_{aw},\qquad
\begin{aligned}
j^\mu_{\Pi,A}&=128f^2\sum_{w\in W}\Im(\bar\phi_wD^\mu\phi_w),\\
j^\mu_{\Pi,\rm em}&=128f^2\sum_{a\in A}\Im(\bar\phi_aD^\mu\phi_a),
\end{aligned}
\]
with $\phi_w=x_{2,w}+ix_{3,w}$ and $\phi_a=x_{a,13}+ix_{a,7}$. Since $[Q,\Pi_A]=0$ for both charges, no $A_\mu A^\mu$ term arises at $x=0$, and the criterion of
\cref{cor:massless} holds at $\Pi_A$ and $\Pi_W$ for $\uu1_A$ and for $\uu1_{\rm em}$. On the Clifford spheres of \cref{p1:prop:spheres}(iii) the formulas are exact,
with $\mathcal L_\Pi=64f^2\,Dn\cdot Dn$ (p1-31, p1-32).
\end{theorem}

\begin{proof}
$D\Pi=[DX,\Pi_A]+O(x^2)$, and $j_\Pi$ at second order is $-f^2\Tr([[Q,X],\Pi_A][DX,\Pi_A])$. p1-31 evaluates both bilinear forms with $\Tr(T_{aw}T_{a'w'})=128\delta$ and the
pairs of \cref{p1:thm:tangent}. The sign pattern is that of \cref{thm:S4}: $[Q,T_{\rm re}]=-T_{\rm im}$ and $[Q,T_{\rm im}]=T_{\rm re}$.
\end{proof}

\begin{theorem}[Curved extension; fixed]
\label{p1:thm:curved}
\begin{enumerate}[label=(\roman*),leftmargin=2.2em]
\item Let $\Omega=\Omega_{\rm Lorentz}+A^aT_a$ with $t_a$ spanning a subalgebra $\mathfrak s\subset\mathfrak k$. Then $D\Pi_P=0$ and $[\mathcal R,\Pi_P]=0$ hold for arbitrary spacetime and
gauge curvature iff $\mathfrak s\subset\so P\oplus\so{P^\perp\cap V_+}\oplus\so6_C$. At $\Pi_A$ and $\Pi_W$ this includes $\uu1_A$, $\uu1_{\rm em}$, $\mathfrak h_{\rm PS}$ and
$\mathfrak h_{\rm SM}$ (p1-35).
\item At every point of $F_3$, if $D_\mu\Pi=[Y_\mu,\Pi]$ with $Y_\mu\in\mathfrak k$, the adapted connection of \cref{prop:adapted} is
$\mathcal C_\mu=[\Pi,D_\mu\Pi]=-\tfrac12(Y_\mu-\mathrm{vol}_PY_\mu\mathrm{vol}_P)$. This is a compact bivector of $\Lambda^2V_+$, so the adapted connection is spin-valued across
the whole Grassmannian family (p1-33).
\end{enumerate}
\end{theorem}

\begin{proof}
(i) The six interior bivectors commute with every $\Pi_P$ and with $\mathfrak k$. So $D\Pi_P=A^a[t_a,\Pi_P]$ and $\mathcal R=\mathcal R_{\rm Lorentz}+F^aT_a$, and both vanish
for arbitrary $A$ and $F$ iff each $t_a$ is in the stabiliser of \cref{p1:thm:classify}. (ii) With $v=\mathrm{vol}_P$ and $v^2=\id$,
$[\Pi,[Y,\Pi]]=\tfrac14[v,[Y,v]]=\tfrac12(vYv-Y)$. $\Ad(v)$ acts on vectors as $-1$ on $P$ and $+1$ on $P^\perp\cap V_+$ and $V_-$. So it maps $\Lambda^2V_+$ and
$\Lambda^2V_-$ to themselves, and the $\Lambda^2V_-$ part of $Y$ cancels.
\end{proof}

\begin{proposition}[Invariant couplings are constant; fixed]
\label{p1:prop:K}
Every $K$-invariant function on $F_3$ is constant. In particular the full alignment energy satisfies $\tfrac1{512}\sum\|[e_ie_j,\Pi]\|^2=16$ on $F_3$, with the sum over the $43$
compact blades. Every spectrum-only potential is constant on $F_3$. For a subalgebra $\mathfrak h\subsetneq\mathfrak k$ and positive couplings, the zero set of $V_{\mathfrak h}$ on
$F_3$ is the fixed set $\operatorname{Fix}(\mathfrak h)=\{\Pi\in F_3:[\mathfrak h,\Pi]=0\}$ (p1-34, p1-42).
\end{proposition}

\begin{proof}
$F_3$ is one $K$-orbit. The $43$ blades are orthogonal of equal norm in the trace form, which $\Ad(K)$ preserves, so $\sum_a\|[t_a,\Ad(g)\Pi]\|^2=\sum_a\|[\Ad(g^{-1})t_a,\Pi]\|^2$
is a sum over another orthonormal basis and has the same value. The value $16$ is sixteen movers, each with $\|[e_ie_j,\Pi_A]\|^2=512$. $V_{\mathfrak h}$ is a sum of
squares with positive weights.
\end{proof}

A coupling that selects a point of $F_3$ must therefore break $K$, and the natural candidates are the $V_{\mathfrak h}$ of proper subalgebras
\cite{Peskin80,Preskill81}. \Cref{p1:prop:duality} adds that every such $V_{\mathfrak h}$ is also invariant under $D$ and under $\Pi\mapsto\id-\Pi$.

% --- source: photon_selection/p1_families#6 (Fixed sets and the colour condition)
\subsection{Fixed sets and the colour condition}
\label{p1:sec:fixed}

Write $U_A=\operatorname{span}(E_2,E_3)$, $U_{\rm em}=\operatorname{span}(w_1,w_2)$ and $W_0=\operatorname{span}(w_3,w_4)$.

\begin{theorem}[Fixed sets in $F_3$; fixed]
\label{p1:thm:fixed}
\begin{center}
\small
\begin{tabular}{@{}lllc@{}}
\toprule
algebra & fixed set in $F_3$ & dimensions & contains\\
\midrule
$\so6_C$, $\su3_C$ & all of $F_3$ & $16$ & all points\\
$\uu1_A$ & $\{P\supset U_A\}\sqcup\{P\perp U_A\}\cong\widetilde{\mathrm{Gr}}_2(\R^6)\sqcup\widetilde{\mathrm{Gr}}_4(\R^6)$ & $8,\ 8$ & $\Pi_{\pm A}$; $\Pi_{\pm W}$\\
$\uu1_{\rm em}$ & $\{P\supset U_{\rm em}\}\sqcup\{P\perp U_{\rm em}\}\cong\widetilde{\mathrm{Gr}}_2(\R^6)\sqcup\widetilde{\mathrm{Gr}}_4(\R^6)$ & $8,\ 8$ & $\Pi_{\pm W}$; $\Pi_{\pm A}$\\
$\R t_Y$ & $\{\Pi_{\pm A}\}\sqcup\{\Pi_{\pm W}\}\sqcup\{A'\oplus W'\}$ & $0,\ 0,\ 6$ & \\
$\su2_L$, $\mathfrak h_{\rm SM}$, $\mathfrak h_{\rm PS}$ & $\{\Pi_A,\Pi_{-A},\Pi_W,\Pi_{-W}\}$ & $0$ & four isolated points\\
$\mathfrak k$ & $\emptyset$ & & \\
\bottomrule
\end{tabular}
\end{center}
In the $t_Y$ row, $A'$ is a two-plane in $A$ and $W'$ a $J_Y$-complex line in $W$. That component is $\widetilde{\mathrm{Gr}}_2(A)\times\C P^1$, connected (p1-36 to p1-39).
\end{theorem}

\begin{proof}
For a connected group, $\Pi_P$ is fixed iff $P$ is invariant under $\rho_+(\mathfrak h)$ (\cref{p1:lem:transpose}). Let $j$ be a complex structure on a subspace $U$ that kills
$U^\perp$. If $P$ is $j$-invariant and $v=v_U+v_\perp\in P$, then $j^2v=-v_U\in P$, so $P=(P\cap U)\oplus(P\cap U^\perp)$ with $P\cap U$ itself $j$-invariant. For $U=U_A$ or
$U_{\rm em}$, $P\cap U\in\{0,U\}$, which gives the two Grassmannians of $U^\perp\cong\R^6$, each connected of dimension $8$. For $U=W$ and $j=J_Y$, $P\cap W$ is a complex
subspace of dimension $0$, $2$ or $4$. For $\su2_L$: $A$ is trivial and $W$ is irreducible without invariant vectors (p1-37). So $P\cap W\in\{0,W\}$, and if $P\cap W=0$
the $W$-component of every vector of $P$ is invariant, hence zero; thus $P\in\{A,W\}$. $\mathfrak h_{\rm SM}$ and $\mathfrak h_{\rm PS}$ contain $\su2_L$ and fix $A$ and $W$. $\so8$
is irreducible on $V_+$. The tangent kernels in p1-36, p1-38 and p1-39 agree with these dimensions, and p1-38 shows that the four points are isolated.
\end{proof}

\begin{proposition}[Unbroken algebras across $F_3$; fixed]
\label{p1:prop:strata}
The stabiliser of $\Pi_P$ in $\mathfrak h_{\rm PS}$ is $\so6_C\oplus\so{P\cap W}\oplus\so{P^\perp\cap W}$. In $\mathfrak h_{\rm SM}$ it is
$\su3_C\oplus\bigl(\mathfrak u_W\cap(\so{P\cap W}\oplus\so{P^\perp\cap W})\bigr)$, with $\mathfrak u_W=\rho(\su2_L\oplus\R t_Y)\cong\uu2$. At random points the dimensions in
$(\mathfrak h_{\rm SM},\mathfrak h_{\rm PS})$ are $(9,16)$ on both $\uu1_{\rm em}$-fixed components, namely $\su3_C\oplus\uu1_{\rm em}$ and $\so6_C\oplus\uu1_{\rm em}$. They are
$(8,16)$ on both $\uu1_A$-fixed components, $(8,15)$ at generic points of $F_3$ and $(12,21)$ at $\Pi_{\pm A},\Pi_{\pm W}$. On the six-dimensional
$t_Y$-component the stabiliser in $\mathfrak h_{\rm SM}$ has dimension $10$ (p1-39, p1-40).
\end{proposition}

\begin{proof}
If $X\in\so W$ preserves $P$ and kills $A$, then $X(P)\subset P\cap W$ and $X(P^\perp)\subset P^\perp\cap W$. So $X$ has image in $(P\cap W)\oplus(P^\perp\cap W)$, and by
antisymmetry it lies in $\so{P\cap W}\oplus\so{P^\perp\cap W}$; the converse is immediate. On $\{P\perp U_{\rm em}\}$ one generically has $P\cap W=0$ and
$P^\perp\cap W=U_{\rm em}$, and $\so{U_{\rm em}}=\R\rho(\tau_3)=\R\rho(B_{\rm em})\subset\mathfrak u_W$. p1-40 computes the dimensions of the intersections and of the stabilisers.
\end{proof}

The $\uu1_{\rm em}$-fixed component through $\Pi_A$ is $\widetilde{\mathrm{Gr}}_4(A\oplus W_0)$. Its tangent at $\Pi_A$ is $A\otimes W_0$, the neutral lower components of the four
doublets of \cref{p1:thm:tangent}. Along it $\su3\oplus\su2_L\oplus\uu1_Y$ breaks to $\su3\oplus\uu1_{\rm em}$, the pattern of a neutral doublet value (interpretive).

\begin{corollary}[Critical points and Hessians; fixed]
\label{p1:cor:critical}
Every smooth function on $F_3$ invariant under the connected group of $\mathfrak h\in\{\su2_L,\mathfrak h_{\rm SM},\mathfrak h_{\rm PS}\}$ has $\Pi_{\pm A}$ and $\Pi_{\pm W}$ as
critical points. Its Hessian there lies in the space of invariant symmetric forms on $A\otimes W$. For $\mathfrak h_{\rm PS}$ this space has dimension $10$, the forms
$S\otimes\delta_W$ with $S$ symmetric on $A$. For $\mathfrak h_{\rm SM}$ it has dimension $16$, the forms $S\otimes\delta_W+N\otimes\omega_Y$ with $N$ antisymmetric and
$\omega_Y(w,w')=\langle w,J_Yw'\rangle$ (p1-41).
\end{corollary}

\begin{proof}
At a fixed point the differential is an invariant covector, and by p1-38 there is none other than zero \cite{Michel71,Palais79}. The Hessian at a critical point is invariant. The
counts follow from Schur's lemma, with the commutant $\R$ of $\so4_W$ on $W$ and the commutant $\C=\operatorname{span}(\id,J_Y)$ of $\mathfrak u_W$; p1-41 computes them.
\end{proof}

\begin{remark}[The fourth rung, in dimensions; fixed]
\label{p1:rem:n4}
At $n=4$ the internal signature is $(11,9)$. The stabiliser of $E_2E_3E_8E_9$ in $\so{11,9}$ has dimension $126$ with compact part $63$, the $n=4$ row of
\cref{prop:towermatter}. The family $F_4=\widetilde{\mathrm{Gr}}_4(\R^{11})$ has dimension $91-63=28$ (p1-43). All $36$ bivectors of $V_-$ commute with every point, so any
colour algebra in $\so{V_-}$ passes the colour test everywhere. At $\Pi_A$ a $\uu1$ on a plane inside $A$ fixes a component of dimension $14=\dim\widetilde{\mathrm{Gr}}_2(\R^9)$,
and one on a plane orthogonal to $A$ fixes one of dimension $20=\dim\widetilde{\mathrm{Gr}}_4(\R^9)$. $\so4_W$ fixes a component of dimension $12=\dim\widetilde{\mathrm{Gr}}_4(\R^7)$
through $\Pi_A$, so at $n=4$ the Pati--Salam fixed points are no longer isolated (p1-44). In general $\dim F_n=4(3n-5)$, and a $\uu1$ on one positive plane fixes components of
dimensions $2(3n-5)$ and $4(3n-7)$: $(2)$, $(8,8)$, $(14,20)$ at $n=2,3,4$ (p1-45).
\end{remark}

% --- source: photon_selection/p1_families#7 (The answer)
\subsection{The answer}
\label{p1:sec:answer}

\begin{enumerate}[label=(\alph*),leftmargin=2.2em]
\item \emph{The family.} The $n\ge3$ analogue of the compact family is $F_n=\{\Pi_P:P\in\widetilde{\mathrm{Gr}}_4(V_+)\}$, the $K_n$-orbit of $\Pi_A$; at $n=2$ it is the paper's $S^4$
(fixed). At $n=3$ it is $\widetilde{\mathrm{Gr}}_4(\R^8)$, of dimension $16$, a compact symmetric space \cite{Helgason}. Its metric is $64$ times the trace metric of the
Grassmannian (\cref{p1:lem:vector}). The paper's $S^4$ is one of the inner Clifford spheres in it, and is not $K$-invariant (fixed).
\item \emph{Colour.} Every point of $F_3$ commutes with all of $\so6_C\supset\su3_C$, and so does every point of $F_n$ with $\so{V_-}$. The other involutive four-blade types never
commute with $\su3_C$ (fixed).
\item \emph{The vacuum.} $\Pi_A$ is the analogue of \cref{def:towervac}. It is Lorentz invariant, neutral under $\uu1_A$ and $\uu1_{\rm em}$, fixed by $\mathfrak h_{\rm PS}$ and
$\mathfrak h_{\rm SM}$, moved by $16$ compact bivectors, and parallel under any connection valued in its stabiliser. It is the point $n=c$ of the embedded second-rung
family. The literal point $n=e$ is not neutral at $n=3$ (fixed).
\item \emph{Data for a selecting coupling.} A coupling that distinguishes the points of $F_3$ is not $K$-invariant (obstructed for $K$-invariant couplings). Its zero set, if it is
an alignment energy, is one of the fixed sets of \cref{p1:thm:fixed}. Its Hessian at the four Pati--Salam points lies in the $10$- or $16$-dimensional spaces of
\cref{p1:cor:critical}. Its fluctuations there are four doublets $(\mathbf1,\mathbf2)_{1/2}$ with a family index (fixed). No alignment energy of a subalgebra of $\mathfrak k$
separates $\Pi_A$, $\Pi_{-A}$, $\Pi_W$ and $\Pi_{-W}$; a coupling to the chiral block does, since on the block $D$ exchanges $P_L$ and $P_R$ (fixed). This is the
counterpart on $F_3$ of \cref{prop:alignment} (interpretive).
\item \emph{The charge.} On $P_*$ the construction's charge $\tfrac12J_1$ is $\tfrac12J$, universal, and anomalous together with $\su2_L$ (fixed; restated). The electromagnetic
charge $B_{\rm em}$ also leaves $\Pi_A$ massless and has four charged complex fluctuations. Which charge, which coupling and which scale select the vacuum is not decided
here (open).
\end{enumerate}

\begin{table}[t]
\centering
\footnotesize
\begin{tabular}{@{}>{\raggedright\arraybackslash}p{8.4cm}>{\raggedright\arraybackslash}p{4.3cm}>{\raggedright\arraybackslash}p{2.2cm}@{}}
\toprule
statement & class & checks\\
\midrule
\cref{p1:prop:n2}: second-rung data, $K_2$-orbit $=S^4=\widetilde{\mathrm{Gr}}_4(\R^5)$, $Q$-fixed $S^2$, constant potentials & restated (\crefrange{thm:towervac}{thm:towercurved}, \cref{prop:towermatter}); orbit and fixed set fixed & p1-1 to p1-7\\
\cref{p1:lem:census}: four-blade census, even-grade rule, $\mathrm{vol}_+=-\mathcal Z$ & fixed & p1-8 to p1-10\\
\cref{p1:lem:transpose}: $\mathrm{vol}_P^{\!\top}=\mathrm{vol}_{\theta P}$, recovery of $P$ & fixed & p1-13\\
\cref{p1:lem:colour}: $\su3_C$ irreducible on $V_-$ & fixed & p1-15\\
\cref{p1:thm:classify}: three $K$-orbits, stabilisers, spectra, colour test & fixed; $n=3$ row of \cref{prop:towermatter} restated & p1-11, p1-12, p1-14, p1-16, p1-42\\
\cref{p1:prop:spheres}: Clifford spheres, embedded second-rung family & fixed & p1-19, p1-32, p1-47\\
\cref{p1:prop:duality}: duality $P\mapsto P^\perp$, invariance of mass matrices & fixed & p1-20, p1-46\\
\cref{p1:thm:vacuum}: $\Pi_A$, $\Pi_W$ neutral, $\mathfrak h_{\rm PS}$- and $\mathfrak h_{\rm SM}$-fixed, not rigid; literal embedding fails & fixed & p1-21 to p1-23, p1-35\\
\cref{p1:thm:tangent}: fluctuations, four doublets $(\mathbf1,\mathbf2)_{1/2}$, charges & fixed & p1-24 to p1-26\\
\cref{p1:prop:block}: $\Pi_AP_*=P_L$, $q_A=J$ on the block, $P_*TP_*=0$, anomaly & fixed; (iv) restated (\cref{prop:PSanom}); anomaly condition inherited & p1-17, p1-18, p1-27, p1-28\\
\cref{p1:lem:vector}: vector model, kinetic term, current, $V_{\mathfrak h}$ & fixed & p1-29, p1-30\\
\cref{p1:thm:kinetic}: $\mathcal L_\Pi$ and $j_\Pi$ at $\Pi_A$ and on spheres & fixed & p1-31, p1-32\\
\cref{p1:thm:curved}: parallel vacua, compact adapted connection & fixed & p1-33, p1-35\\
\cref{p1:prop:K}: $K$-invariant couplings constant; zero sets & fixed; selection by $K$-invariant couplings obstructed & p1-34, p1-42\\
\cref{p1:thm:fixed}, \cref{p1:prop:strata}: fixed sets, unbroken algebras & fixed & p1-36 to p1-40\\
\cref{p1:cor:critical}: critical points and Hessians & fixed; criticality principle inherited & p1-38, p1-41\\
\cref{p1:rem:n4}: fourth rung & fixed; $n=4$ row of \cref{prop:towermatter} restated & p1-43 to p1-45\\
family reading of $\so4_A$; electroweak pattern on the neutral component & interpretive & ---\\
which charge and coupling select the vacuum & open & ---\\
\bottomrule
\end{tabular}
\caption{Claims of \cref{p1:sec:main} and their classes. The checks are the \texttt{PASS} lines printed by \texttt{p1\_families.py}.}
\label{p1:tab:claims}
\end{table}

% --- source: photon_selection/p2_alignment#0 ((head))
% No new references beyond those of p1 and p3 (Coleman--Weinberg, Peskin80, Preskill81, Witten83 are cited in p3).

% --- source: photon_selection/p2_alignment#1 (Alignment by a gauged subalgebra)
\section{Alignment by a gauged subalgebra}
\label{p2:sec:main}

\Cref{p1:sec:main} found the compact family $F_3$ of vacuum projectors at the third rung, and showed that couplings invariant under the whole
compact internal algebra $\mathfrak k$ are constant on it. This section asks what a gauged subalgebra $\mathfrak h\subset\mathfrak k$ does. The answers are
as follows. The gauge field's mass matrix at a vacuum $\Pi$ has trace $V_{\mathfrak h}(\Pi)=\sum_ag_a^2\|[t_a,\Pi]\|^2$, its kernel is the unbroken
subalgebra (fixed). Its zero set on $F_3$ is the set of vacua that leave $\mathfrak h$ unbroken (fixed). Which point is the vacuum is decided by the sign of the
coefficient with which $V_{\mathfrak h}$ enters the vacuum energy. That sign is the hypothesis (H3) of \cref{p3:thm:selection}; the one-loop gauge contribution
supplies it only in a cutoff scheme (inherited; scheme dependent). Identifiers p2-$n$ refer to the \texttt{PASS} lines of
\texttt{p2\_alignment.py}.

\paragraph{Conventions.} As in \cref{p1:sec:main,p3:sec:main}. On $\mathfrak k$ the invariant form here is $\langle X,Y\rangle=-\Tr(XY)/128$, in which the unit-rate
rotations $\tfrac12e_ie_j$ are orthonormal (p2-1). This differs by a factor $4$ from the form $-\Tr(XY)/512$ of \cref{p3:sec:main}; norms below are quoted in the
form of this section.

\begin{proposition}[The alignment energy is the trace of the mass matrix; fixed]
\label{p2:prop:mass}
For $\mathfrak h$ with orthonormal basis $t_a$ and couplings $g_a>0$, put $M^2_{ab}(\Pi)=g_ag_b\Tr([t_a,\Pi][t_b,\Pi])$. Then $M^2$ is symmetric and positive
semidefinite, $\Tr M^2=V_{\mathfrak h}(\Pi)$, $\dim\ker M^2$ is the dimension of the stabiliser of $\Pi$ in $\mathfrak h$, and $V_{\mathfrak h}$ is invariant under the
group generated by $\mathfrak h$ (p2-2, p2-3). On $F_3$, $V_{\mathfrak h}(\Pi_P)=64\sum_ag_a^2\|[\rho_+t_a,p_P]\|^2$ is computed on $V_+=\R^8$, and the colour couplings drop
out (p2-4). For $\mathfrak h=\uu1_A,\uu1_{\rm em},\mathfrak h_{\rm SM},\mathfrak h_{\rm PS},\dots$ it has the closed forms of p2-13 in the principal angles between $P$ and the
planes $U_A$, $U_{\rm em}$ and $W$.
\end{proposition}

\begin{proposition}[Gauging all of $\mathfrak k$ selects nothing; fixed]
\label{p2:prop:allk}
With $\mathfrak h=\mathfrak k$ and $\Ad(K)$-invariant couplings, the mass matrix has spectrum $\{128g_8^2\ (\times16),\ 0\ (\times27)\}$ at every point of $F_3$, so every
spectral function of $M^2$, the full one-loop potential included, is constant on $F_3$ (p2-6).
\end{proposition}

So the selection of a point needs a proper subalgebra, as \cref{p1:prop:K} already showed for $K$-invariant couplings in general.

\begin{proposition}[Zero sets and Goldstone counts; fixed]
\label{p2:prop:zero}
For each $\mathfrak h$ the global minimum of $V_{\mathfrak h}\ge0$ on $F_3$ is its zero set, the $\mathfrak h$-fixed points of \cref{p1:thm:fixed}: for $\uu1_A$ and $\uu1_{\rm em}$ (with or
without $\su3_C$) the two components $P\supset U$ and $P\perp U$, for $\mathfrak h_{\rm SM}$, $\mathfrak h_{\rm PS}$ and $\mathfrak h_{\rm PS}\oplus\uu1_A$ the four points
$\Pi_{\pm A},\Pi_{\pm W}$ (p2-16: $60$ random starts per algebra, random couplings, all reaching $V_{\mathfrak h}<10^{-12}$ relative). At $\Pi_A$ the Hessian of $V_{\mathfrak h}$ is
positive semidefinite with kernel $[\mathfrak c_{\mathfrak k}(\mathfrak h),\Pi_A]$: $\uu1_A$ and $\uu1_{\rm em}$ leave eight Goldstone directions, and every $\mathfrak h$
containing $\su2_L$ or $\so4_W$ leaves none (p2-14, p2-15, p2-18).
\end{proposition}

\begin{proposition}[A second critical set of $\mathfrak h_{\rm SM}$; fixed]
\label{p2:prop:metastable}
$P_0=\operatorname{span}(E_2,E_3,E_7,E_{13})$ is a critical point of $V_{\mathfrak h_{\rm SM}}$, on a critical set $\widetilde{\mathrm{Gr}}_2(A)\times\C P^1$. Relative to
$\Tr(d\Pi^2)$ its Hessian has eigenvalues $0$ ($\times6$, along the critical set), $\tfrac12(g_Y^2-g_2^2)$ ($\times8$) and $2g_Y^2$ ($\times2$). So it is a saddle for
$g_Y<g_2$ and a local, non-global minimum for $g_Y>g_2$. There the unbroken algebra is $\su3_C\oplus\uu1_{T^3_L}\oplus\uu1_Y$, of dimension $10$ and containing
$B_{\rm em}$, and only the $W^\pm$ are massive (p2-17, p2-20).
\end{proposition}

This metastable point is not electroweak breaking: the $Z$ stays massless, and it exists only for $g_Y>g_2$, a ratio of couplings that is
not physical at low energies and that the equal-coupling convention of \cref{p3:thm:photon}(d) excludes (interpretive comparison).

\begin{remark}[The sign hypothesis (H3); inherited, scheme dependent]
\label{p2:rem:sign}
In the one-loop effective potential with a Euclidean momentum cutoff $\Lambda$, a real degree of freedom of mass $m$ contributes
$\tfrac{\Lambda^2m^2}{32\pi^2}+\tfrac{m^4}{64\pi^2}\bigl(\log\tfrac{m^2}{\Lambda^2}-\tfrac12\bigr)+O(m^6/\Lambda^2)$, and a massive vector in Landau gauge counts three (p2-5,
inherited \cite{ColemanWeinberg73}). The leading term is $\tfrac{3\Lambda^2}{32\pi^2}\Tr M^2=\tfrac{3\Lambda^2}{32\pi^2}V_{\mathfrak h}$, so in this scheme $\xi_1>0$ at order
$\Lambda^2$. The coefficient is regulator dependent (it is absent in dimensional regularisation), the $m^4\log m^2$ terms are not controlled, and the fermion
loops of the matter enter with the opposite sign. The sign is proved for vector-like theories \cite{Witten83}. Here (H3) therefore remains a hypothesis;
the curvature masses of $M^2$ are not pole masses, and no gauge-fixing or scheme dependence of the vacuum energy has been examined.
\end{remark}

\begin{remark}[The second rung; fixed]
\label{p2:rem:n2}
At $n=2$, with $\mathfrak h=\uu1_Q\oplus\so3_{cde}$ on the family $S^4$, the global minimum is the uncharged two-sphere (photon massless, $\so3\to\so2$) for
$g_3^2<\tfrac12g_Q^2$, and the charged circle (photon massive, $\so3$ unbroken) for $g_3^2>\tfrac12g_Q^2$; at $g_3^2=\tfrac12g_Q^2$ the potential is constant
(p2-8 to p2-10). Which charge is massless is decided by the ratio of the couplings, not by the algebra.
\end{remark}

\paragraph{What this section establishes.} Gauging a proper subalgebra with a positive alignment coefficient selects the points that leave it unbroken
(fixed). The positivity of the coefficient is a hypothesis with one-loop support in a cutoff scheme only (inherited). The vacua selected by $\mathfrak h_{\rm SM}$ are
the electroweak-symmetric points $\Pi_{\pm A},\Pi_{\pm W}$, and the only other critical set of $V_{\mathfrak h_{\rm SM}}$ breaks $\su2_L$ to $\uu1_{T^3_L}$, which is not
the observed pattern. Electroweak breaking is not produced by alignment, in agreement with \cref{p3:prop:neutral}.

% --- source: photon_selection/p3_identification#0 ((head))
% New references: sections/bib_p3.tex.

% --- source: photon_selection/p3_identification#1 (Which charge is the photon, and the residual choice)
\section{Which charge is the photon, and the residual choice}
\label{p3:sec:main}

The second-rung construction gauges the charge $Q=\tfrac12J_1$ and carries it, in the curved extension, by the normal $SO(2)$ connection
\cref{thm:S4,thm:towercurved}. At the third rung that charge is $Q_A=\tfrac12q_A=\tfrac12E_2E_3$, and the matter block of \cref{sec:pstar} carries a
second candidate, the electric charge $B_{\rm em}$. This section decides which of the two is the photon on the family $F_3$ of \cref{p1:sec:main}, and what is
left to choose once a gauged subalgebra has aligned the vacuum. The results are as follows. The charge operator of $Q_A$ is half the involution $\mathcal U$ that cuts
out $P_*$, so on the block $Q_A$ equals $\tfrac12$ on every state: it is half the fermion number, and it is anomalous together with $\su2_L$ (fixed). It is not
electromagnetism on that matter, whatever the vacuum (obstructed). The electric charge leaves the four aligned vacua massless, has four charged complex
fluctuations with the electron's charge magnitude, and acts on the block with the Standard-Model charges and without anomaly. With $B_{\rm em}$ gauged, the
massless vector of the projector current is the candidate for the electromagnetic photon of the block (charge identification fixed; Maxwell term adopted;
an electroweak stage, which the potential of \cref{p3:sec:selection} disfavours, is still missing), and the normal $U(1)$ is a separate abelian field (fixed). Of the three
counts against the normal $U(1)$ in \cref{sec:assessment} only the sign count changes status, and a fourth count, the charge on the block, decides.
After alignment by $\mathfrak h_{\rm SM}$ four vacua remain. A Klein group of symmetries of the projector sector permutes them simply transitively, so no
alignment energy, kinetic term or current separates them (fixed). The overlaps with $P_*$, with the lepton projector and with the chirality halves do
not separate them either (fixed). The overlap with the half $P_L$ of the block on which $\su2_L$ acts selects the pair $\{\Pi_A,\Pi_{-W}\}$, and the two
members of the pair coincide on the half of the module that contains the block (fixed). A single point needs a coupling that sees the other half, for
instance the overlap with $\Pfam$ (constructed). Identifiers p3-$n$ refer to the \texttt{PASS} lines of \texttt{p3\_identification.py}
(\cref{p3:tab:claims}).

\paragraph{Conventions.} Definitions and notation are those of \cref{p1:sec:main}; in particular $F_3$, $\Pi_{\pm A}$, $\Pi_{\pm W}$, $Q_A$, $B_{\rm em}$, $t^{L,R}_i$,
$t_Y$, $H_C$, $\rho_+$ and $\mathrm{vol}_+$ are as there. For a compact generator $X$ commuting with $J$ the charge operator is $\hat X:=-JX$. It is symmetric, and
$X=J\hat X$ acts as $i\hat X$ in the complex structure $J$ \cref{sec:pstar}. The inner product on $\mathfrak k$ is $\langle X,Y\rangle=-\Tr(XY)/512$. For it the unit
blades $e_ie_j$ have norm one, $|t^L_3|^2=\tfrac18$ and $|t_Y|^2=\tfrac5{24}$. Orthonormal bases are $2\sqrt2\,t^L_i$, $\sqrt{24/5}\,t_Y$ and an orthonormal basis of
$\su3_C$ for $\mathfrak h_{\rm SM}$, with couplings $g_2$, $g_Y$, $g_3$, and the unit blades of $\so4_W\oplus\so6_C$ for $\mathfrak h_{\rm PS}$, with one coupling
$g$. We write $U_A=\operatorname{span}(E_2,E_3)$, $A_0=\operatorname{span}(E_8,E_9)$, $U_{\rm em}=\operatorname{span}(E_7,E_{13})$, $W_0=\operatorname{span}(E_{14},E_{15})$
and $P_{L,R}=P_*\tfrac12(\id\mp\Gamma)$. For an oriented four-plane $P\subset V_+$ with oriented orthonormal frame $(f_1,\dots,f_4)$, written as a $4\times8$ matrix $F$ in
the generator basis of $V_+$, the Plücker coordinates are $p_I(P)=\det F_I$ for the $70$ four-subsets $I$ of that basis, and $e_I$ is the blade of $I$ in increasing
order. We write $p_A$ and $p_W$ for $I=A$ and $I=W$. The family projector is $\Pfam=\tfrac18\prod_i(\id+X_i)$ \cref{sec:projectors}, and
$Q_v=P_L+\gamma^0P_L\gamma^0$ \cref{prop:alignment}.

\begin{definition}[The lepton projector]
\label{p3:def:plep}
$P_{\rm lep}$ is the spectral projector of the colour charge $Q_c=-JH_C$ for the eigenvalue $-3$. With $s_k=J\,E_{9+k}E_{15+k}$, $k=1,2,3$, one has
$Q_c=s_1+s_2+s_3$ and $P_{\rm lep}=\tfrac18\prod_k(\id-s_k)$.
\end{definition}

On $P_*$ the eigenvalue $-3$ of $Q_c$ is the colour singlet \cref{prop:PSledger}. \Cref{p3:prop:restated} shows that $P_{\rm lep}$ reproduces the row of
\cref{prop:alignment}, where the colour-singlet lepton projector is named but not written out.

% --- source: photon_selection/p3_identification#2 (Two charges on the matter)
\subsection{Two charges on the matter}
\label{p3:sec:charges}

\begin{proposition}[The construction's charge is half of $\mathcal U$; fixed, and restated]
\label{p3:prop:charges}
\begin{enumerate}[label=(\roman*),leftmargin=2.2em]
\item $\hat Q_A=-JQ_A=\tfrac12\mathcal U$ on $\R^{512}$, where $\mathcal U=-Jq_A$ is the third involution of $P_*$. Since $\mathcal UP_*=P_*$ and $q_A$ commutes with $\mathcal X$, $\mathcal Z$ and $\mathcal U$,
\[
q_AP_*=JP_*,\qquad \hat Q_AP_*=\tfrac12P_* .
\]
The first identity is the statement $q_A|_{V_*}=J$ of \cref{sec:pstar} (p3-1).
\item On $V_*=\operatorname{im}P_*$ the operators $\Gamma$, $\hat T^3_L=-Jt^L_3$, $\hat Y=-Jt_Y$, $\hat Q_{\rm em}=-JB_{\rm em}$, $\hat Q_A$ and $\BmL=-JH_C/3$ commute. Their
joint spectrum is that of \cref{p3:tab:charges} (p3-2).
\item On the Weyl block $\tfrac12(\id-\Gamma)$ of the spectator action, of rank $256$, $\hat Q_A$ takes the values $\pm\tfrac12$, each $128$ times. There
$\hat Q_{\rm em}$ takes the values $\pm1$ ($16$ times each), $\pm\tfrac23$ ($48$), $\pm\tfrac13$ ($48$) and $0$ ($32$), the charges of the two families and two mirror
families of \cref{prop:spectator} (p3-3).
\end{enumerate}
\end{proposition}

\begin{table}[t]
\centering
\small
\begin{tabular}{@{}lcccccccc@{}}
\toprule
field & $\Gamma$ & $\hat T^3_L$ & $\hat Y$ & $\hat Q_{\rm em}$ & $\hat Q_A$ & $\BmL$ & real dim. & complex dim.\\
\midrule
$Q_L$ & $-1$ & $\pm\tfrac12$ & $\tfrac16$ & $\tfrac23,\ -\tfrac13$ & $\tfrac12$ & $\tfrac13$ & $12+12$ & $6+6$\\
$L_L$ & $-1$ & $\pm\tfrac12$ & $-\tfrac12$ & $0,\ -1$ & $\tfrac12$ & $-1$ & $4+4$ & $2+2$\\
$u_R$ & $+1$ & $0$ & $\tfrac23$ & $\tfrac23$ & $\tfrac12$ & $\tfrac13$ & $12$ & $6$\\
$d_R$ & $+1$ & $0$ & $-\tfrac13$ & $-\tfrac13$ & $\tfrac12$ & $\tfrac13$ & $12$ & $6$\\
$e_R$ & $+1$ & $0$ & $-1$ & $-1$ & $\tfrac12$ & $-1$ & $4$ & $2$\\
$\nu_R$ & $+1$ & $0$ & $0$ & $0$ & $\tfrac12$ & $-1$ & $4$ & $2$\\
\bottomrule
\end{tabular}
\caption{Joint eigenvalues on $V_*$ (p3-2). Complex dimensions refer to the complex structure $J$; one internal Weyl component has four real, or two complex,
dimensions. The columns $\hat Y$ and $\hat Q_{\rm em}$ restate \cref{prop:PSledger}. The column $\hat Q_A$ is constant.}
\label{p3:tab:charges}
\end{table}

\begin{proof}
(i) $\hat Q_A=-J\tfrac12q_A=\tfrac12\mathcal U$. As index sets, $\mathcal X\sim I\cup A\cup C$ with $I$ the interior, $\mathcal Z\sim A\cup W$ and $\mathcal U\sim\{8,9\}\cup W\cup C$. The bivector $q_A$ shares two
indices with $\mathcal X$ and with $\mathcal Z$ and none with $\mathcal U$, so by equation~\cref{eq:bladecomm} it commutes with all three. $\mathcal X$, $\mathcal Z$, $\mathcal U$ commute, so $\mathcal UP_*=P_*$, that is
$-Jq_AP_*=P_*$; multiplying by $J$ and using $J^2=-\id$ gives $q_AP_*=JP_*$. (ii) The operators are the charge operators of commuting compact elements, together with
$\Gamma$, which commutes with every internal bivector; p3-2 diagonalises them jointly on $V_*$. (iii) $\mathcal U$ commutes with $\Gamma$ and is traceless on the Weyl block;
p3-3 computes both spectra there.
\end{proof}

\begin{corollary}[What the construction's photon couples to]
\label{p3:cor:notem}
The items are classed as fixed, restated or obstructed.
\begin{enumerate}[label=(\roman*),leftmargin=2.2em]
\item A gauge field coupled through $Q_A$ couples to every Weyl fermion of the block with the same charge $\tfrac12$, on both chiralities. In the complex structure
$J$ each of the sixteen Weyl fermions has fermion number one (\cref{sec:pstar}). The coupling is therefore to $\tfrac12F$, $F=3B+L$, a universal vector charge
(fixed).
\item No real $c$ gives $\hat Q_AP_*=c\,\hat Q_{\rm em}P_*$, since $\hat Q_{\rm em}$ vanishes on the neutrinos and $\hat Q_A$ does not (p3-4). The statement
involves only $P_*$, so the construction's charge is not the electric charge on the block for any vacuum projector (obstructed).
\item $\hat Q_{\rm em}$ has the same spectrum on $P_L$ as on $P_R$, so its cubic, gravitational and mixed colour anomalies cancel between the chiralities, and
$\Tr(P_L\hat Q_{\rm em}(\hat T^3_L)^2)=0$ (fixed; the anomaly conditions are inherited \cite{BouchiatIliopoulosMeyer72,GeorgiGlashow72}). For $\hat Q_A$ the same
traces vanish except $\Tr(P_L\hat Q_A(\hat T^3_L)^2)=4$. This is the mixed $SU(2)_L^2\,U(1)_A$ anomaly of \cref{prop:PSanom} at $a=\tfrac12$, $b=0$ (restated;
p1-28, p3-5). With $P_*$ as the matter and $\su2_L$ gauged, $\uu1_A$ cannot be gauged (obstructed under these hypotheses).
\end{enumerate}
\end{corollary}

\begin{proof}
(i) and (ii) read off \cref{p3:tab:charges}. (iii) is p3-5. The anomaly of $\hat Q_A$ is that of fermion number in the Standard Model, which $SU(2)_L$ instantons
violate \cite{tHooft76}.
\end{proof}

The universality of $Q_A$ on matter is therefore not a property of the vacuum. It is the factor $\tfrac12(\id+\mathcal U)$ of $P_*$: the block is cut out as an eigenspace of
the construction's charge (\cref{p3:prop:charges}(i)), and on the rest of the spectator Weyl block the same charge takes the value $-\tfrac12$
(\cref{p3:prop:charges}(iii)). At the second rung the coupled model and the vacuum share the charge $\tfrac12J_1$ \cref{sec:photon}. At the third rung that
shared charge is half the fermion number.

\begin{remark}[What would make the normal connection admissible; fixed, with options open]
\label{p3:rem:anomaly}
The normal $SO(2)$ connection of the curved extension \cref{thm:towercurved} is the component of the connection along $Q_A$ (\cref{p3:thm:photon}(c)).
With $P_*$ as the matter and $\su2_L$ gauged it is anomalous, so before it can be treated as a physical gauge field it must be excluded, cancelled or
decoupled (PH-05). \emph{Cancelled}: on the whole spectator Weyl block $\tfrac12(\id-\Gamma)$ the mixed trace is $0$, because the rest of the block, two mirror
families with $\hat Q_A=-\tfrac12$, contributes $-4$ (z2-2); cancellation therefore needs the mirror matter that $P_*$ removes. \emph{Excluded}: the
connection is kept flat or non-dynamical; the curved extension then loses its normal gauge field. \emph{Decoupled}: a mass for it, for instance by a
Green--Schwarz or St\"uckelberg mechanism, which is not constructed. None of the three is established here, and the normal field is not treated as physical.
\end{remark}

% --- source: photon_selection/p3_identification#3 (The counts against the normal U(1)U(1), and the dichotomy)
\subsection{The counts against the normal \texorpdfstring{$U(1)$}{U(1)}, and the dichotomy}
\label{p3:sec:counts}

\Cref{p3:tab:counts} takes the three counts of \cref{sec:assessment} to the third rung and adds the count of \cref{p3:cor:notem}. The four aligned vacua
$\Pi_0\in\{\Pi_{\pm A},\Pi_{\pm W}\}$ commute with both $Q_A$ and $B_{\rm em}$ (\cref{p1:thm:vacuum}), so both charges can be compared at the same vacuum.

\begin{table}[t]
\centering
\footnotesize
\begin{tabular}{@{}>{\raggedright\arraybackslash}p{3.1cm}>{\raggedright\arraybackslash}p{4.2cm}>{\raggedright\arraybackslash}p{4.4cm}>{\raggedright\arraybackslash}p{2.6cm}@{}}
\toprule
count & normal $U(1)$, generator $Q_A$ & $\uu1_{\rm em}$, generator $B_{\rm em}$ & change\\
\midrule
(i) a spinless test particle cannot see it \cref{prop:constraints} and \cref{thm:wlgauge}
& holds: it concerns the metric charge of the normal plane; the field couples through the spinor sector, at $\Pi_0$ to four complex scalars $\phi_w$ and to the block
& not applicable, since $B_{\rm em}$ is not a metric charge; it couples through the spinor sector to four complex scalars $\phi_a$ and to the block, and a neutral spinless field does not see the ordinary photon either
& none in substance (restated; interpretive reading)\\
(ii) the adapted generator does not commute with the physical generators \cref{sec:transport}
& removed: $[Q_A,\Pi_0]=0$, $D\Pi_0=0$, and the adapted correction is compact and commutes with the interior gammas \cref{thm:towercurved}, \cref{p1:thm:curved}
& removed in the same way
& none at $n=3$; removed already at $n=2$ (fixed)\\
(iii) it gravitates with the wrong sign \cref{cor:ein6} and \cref{thm:ghostvec}
& holds for the Kaluza--Klein vector of the normal plane
& not covered: the Maxwell term is adopted, and a positive one exists since $\langle B_{\rm em},B_{\rm em}\rangle>0$; read as a Kaluza--Klein vector, the plane $U_{\rm em}$ has the sign of time, like the normal plane
& from forced to adopted (constructed); the Kaluza--Klein reading is interpretive\\
(iv) the charge on the block, \cref{p3:cor:notem}
& universal $\tfrac12$, anomalous with $\su2_L$
& the Standard-Model charges, anomaly free
& decides (obstructed for $Q_A$, fixed for $B_{\rm em}$)\\
\bottomrule
\end{tabular}
\caption{The counts against the normal $U(1)$ at the third rung (p3-4, p3-5). In row (iii), $E_2,E_3,E_7,E_{13}$ square to $+\id$, as the time generator $\gamma^0=E_1$
does, and the colour generators to $-\id$ (p3-4).}
\label{p3:tab:counts}
\end{table}

Row (iii) rests on \cref{cor:ein6} and \cref{thm:ghostvec}: there a vector obtained by reduction on a plane of the sign of time has a kinetic term of the wrong sign
\cite{Hull98}. That theorem is about a metric origin of the field. An internal gauge field with an adopted Maxwell term is not covered by it. For a compact generator
the invariant form is positive (\cref{thm:sig}, \cref{p3:tab:counts}), so a Maxwell term of the healthy sign exists, as in ordinary Yang--Mills theory. If $\uu1_{\rm em}$
were itself obtained by reduction, its weak part would rotate the plane $U_{\rm em}$, of the sign of time, and the same mechanism would apply to that part and not to the
colour part $\tfrac16H_C$ (interpretive). The same holds for $Q_A$, whose plane $U_A$ also has the sign of time. Count (iii) therefore separates the two fields only by
their origin, and not by their generator.

The dichotomy at the end of \cref{sec:transport} reads: either the normal $U(1)$ is not electromagnetic, or the paired-mode projector is not the observation projector.
\cref{sec:photon} resolves it in one direction: if the normal $U(1)$ is to be electromagnetism, the paired-mode projector must go. At the third rung count (iv)
decides the first alternative for every projector, because it is a property of $P_*$ through its factor $\mathcal U$: on the block matter the normal $U(1)$ is not electromagnetic
(obstructed). The replacement of the paired-mode projector is still available, as $\Pi_A$ or any point of the aligned set, and the charge that uses it is $B_{\rm em}$.
At $n=3$ both alternatives of the dichotomy hold, and neither is a choice.

% --- source: photon_selection/p3_identification#4 (Identification at the aligned vacua)
\subsection{Identification at the aligned vacua}
\label{p3:sec:ident}

\begin{proposition}[Joint fixed points of the two charges; fixed]
\label{p3:prop:jointfixed}
The points of $F_3$ commuting with both $Q_A$ and $B_{\rm em}$ form six components: the isolated points $\Pi_{\pm(U_A\oplus U_{\rm em})}$, the isolated points $\Pi_{\pm R}$ with
$R=\operatorname{span}(E_8,E_9,E_{14},E_{15})$, and two copies of $\widetilde{\mathrm{Gr}}_2(\R^4)$, of dimension $4$,
\[
\{\Pi_{U_A\oplus P'}:P'\in\widetilde{\mathrm{Gr}}_2(R)\},\qquad\{\Pi_{U_{\rm em}\oplus P'}:P'\in\widetilde{\mathrm{Gr}}_2(R)\}.
\]
$\Pi_{\pm A}$ lie on the first of these, with tangent $A_0\otimes W_0$, and $\Pi_{\pm W}$ on the second (p3-6).
\end{proposition}

\begin{proof}
By \cref{p1:thm:fixed}, $\Pi_P$ commutes with $Q_A$ iff $P\supset U_A$ or $P\perp U_A$, and with $B_{\rm em}$ iff $P\supset U_{\rm em}$ or $P\perp U_{\rm em}$. The four
combinations give $P=U_A\oplus U_{\rm em}$; $P=U_A\oplus P'$ with $P'\subset(U_A\oplus U_{\rm em})^\perp=R$; $P=U_{\rm em}\oplus P'$; and $P=R$. An orientation of a direct sum
is a pair of orientations of the summands up to a common reversal. This gives $\widetilde{\mathrm{Gr}}_2(R)$ in the middle cases and two points in the outer ones. Finally
$A=U_A\oplus A_0$ and $W=U_{\rm em}\oplus W_0$ with $A_0,W_0\subset R$. p3-6 computes the joint tangent kernels $0$, $0$, $4$, $4$ at representatives.
\end{proof}

\begin{proposition}[Charges of the fluctuations; fixed]
\label{p3:prop:torus}
At each of $\Pi_{\pm A}$, $\Pi_{\pm W}$ the sixteen tangent directions decompose under the torus $\uu1_A\times\uu1_{\rm em}$ into complex lines of weights
\[
U_A\otimes U_{\rm em}:\ (1,1)\oplus(1,-1),\qquad U_A\otimes W_0:\ 2\times(1,0),\qquad A_0\otimes U_{\rm em}:\ 2\times(0,1),
\]
and four real neutral directions $A_0\otimes W_0$ (p3-7). Eight real directions are charged under each $\uu1$, four under both and four under neither, in agreement
with \cref{p1:thm:tangent}(iii),(iv).
\end{proposition}

\begin{proof}
$\rho_+(Q_A)$ is the rotation of $U_A$ at unit rate and $\rho_+(B_{\rm em})$ the rotation of $U_{\rm em}$; on $A\otimes W$ the first acts on the first factor and the
second on the second factor.
\end{proof}

\begin{theorem}[The photon of the aligned vacua; fixed, with the gauged generator a choice]
\label{p3:thm:photon}
Let $\Pi_0\in\{\Pi_{\pm A},\Pi_{\pm W}\}$. Let the internal connection carry a field $A_\mu$ along a compact generator $Q\in\mathfrak k$, acting on spinors by $Q$ and on the
projector by $\ad Q$, and adopt a Maxwell term.
\begin{enumerate}[label=(\alph*),leftmargin=2.2em]
\item For $Q=B_{\rm em}$ the four requirements on an electromagnetic photon of the block hold. (1) The generator is the electric charge: $\hat Q_{\rm em}=\hat T^3_L+\hat Y$ on
$V_*$, with the values of \cref{p3:tab:charges}. (2) $[B_{\rm em},\Pi_0]=0$, so the kinetic term gives no $A_\mu A^\mu$ term. (3) Charged fluctuations exist: the four complex
scalars $\phi_a$ of \cref{p1:thm:tangent}(iii), with current $j^\mu_{\Pi,\rm em}=128f^2\sum_a\Im(\bar\phi_aD^\mu\phi_a)$ (\cref{p1:thm:kinetic}). (4) The matter carries the
Standard-Model charges without anomaly (\cref{p3:cor:notem}(iii)). Both currents have the same unit of charge: $B_{\rm em}^2=-\hat Q_{\rm em}^2$ on $V_*$, with value $-1$ on
the electron, and $(\ad B_{\rm em})^2=-1$ on the charged tangent directions (p3-8). The scalars of the projector current carry the electron's charge magnitude, and no
normalisation is free.
\item For $Q=Q_A$ requirements (2) and (3) hold and (1) and (4) fail (\cref{p3:cor:notem}). The massless vector couples to $\tfrac12F$ on the block and to four complex
scalars $\phi_w$ of charge $1$, twice the fermion charge (p3-8). It is not the electromagnetic photon (obstructed).
\item $Q_A$ and $B_{\rm em}$ commute and are orthogonal. $Q_A$ commutes with $\mathfrak h_{\rm SM}$ and $\mathfrak h_{\rm PS}$ and lies in the centraliser $\so4_A$ of
$\mathfrak h_{\rm PS}$ in $\mathfrak k$ (p3-4). The normal $SO(2)$ connection of \cref{thm:towercurved}, whose plane is the surplus pair $E_2,E_3$ of the first
constituent, is the component of the connection along $Q_A$. It is a second abelian field, not the photon (fixed).
\item All twelve generators of $\mathfrak h_{\rm SM}$ commute with $\Pi_0$, so with $\mathfrak h_{\rm SM}$ gauged no gauge boson is massive at $\Pi_0$, and $\uu1_{\rm em}$
is singled out only by an electroweak stage. The projector's own neutral directions provide one. At $\Pi_\theta=\Ad(\exp\tfrac\theta2E_8E_{14})\Pi_A$, $\theta\notin\pi\Z$, the
mass matrix of $\mathfrak h_{\rm SM}$ has the eigenvalue $0$ nine times (gluons and photon), $m_W^2$ twice and $m_Z^2$ once, with
\[
\frac{m_W^2}{m_Z^2}=\frac1{1+\tfrac35\,g_Y^2/g_2^2},
\]
and its massless electroweak combination is $B_{\rm em}$ (p3-9). At $g_Y=g_2$ the ratio is $\tfrac58=1-\tfrac38$, the value $\cos^2\theta_W$ for the
$\sin^2\theta_W=\tfrac38$ of \cref{thm:compat}. That value reflects the normalisation $|t^L_3|^2=\tfrac18$, $|t_Y|^2=\tfrac5{24}$ of the invariant form together with
the assumption $g_Y=g_2$ for the orthonormalised generators (z2-6). It is the familiar unification-scale value. A physical prediction would need the scale at
which the couplings are equal, the particle spectrum, the running down to the electroweak scale and the thresholds, none of which is given (PH-09).
\end{enumerate}
\end{theorem}

\begin{proof}
(a) (1) is \cref{p3:prop:charges}(ii). (2) and (3) are \cref{p1:thm:vacuum,p1:thm:tangent,p1:thm:kinetic}. At $\Pi_{-A}=\id-\Pi_A$ the tangents change sign and nothing
else changes, and at $\Pi_{\pm W}$ \cref{p1:thm:tangent}(v) applies. For the unit: on $V_*$, $B_{\rm em}=J\hat Q_{\rm em}$ with $J$ commuting and $J^2=-\id$, so
$B_{\rm em}^2=-\hat Q_{\rm em}^2$; on the tangent, $[B_{\rm em},T_{a,7}]=T_{a,13}$ and $[B_{\rm em},T_{a,13}]=-T_{a,7}$ (\cref{p1:thm:tangent}). (b) p3-8 gives
$(\ad Q_A)^2=-1$ on $U_A\otimes W$ and $Q_A^2=-\tfrac14$ on $V_*$. (c) is p3-4 together with the form $\Omega=\Omega_{\rm Lorentz}+AQ$, $Q=\tfrac12ab$, of the curved
extension \cref{thm:towercurved}. (d) $\su3_C$ acts trivially on $V_+$. In the vector model of \cref{p1:lem:vector}, $p_\theta=p_{\operatorname{span}(E_2,E_3,E_9)}+uu^{\!\top}$ with
$u=\cos\theta\,E_8\pm\sin\theta\,E_{14}$. For $x\in\rho_+(\mathfrak h_{\rm SM})$, which kills $A$, $[x,uu^{\!\top}]=au^{\!\top}+ua^{\!\top}$ with $a=xu=\pm\sin\theta\,xE_{14}$, so
the mass matrix is $128\sin^2\theta$ times the Gram matrix of the vectors $g_a\rho_+(t_a)E_{14}$. The vectors $\rho_+(t^L_i)E_{14}$ are orthogonal of length $\tfrac12$,
because the $2\rho_+(t^L_i)$ are anticommuting complex structures, and $\rho_+(t_Y)E_{14}=-\rho_+(t^L_3)E_{14}$, because $B_{\rm em}$ kills $E_{14}$. In the basis of the
conventions the Gram matrix is $2g_2^2$ on the first two generators, and on $(2\sqrt2\,t^L_3,\sqrt{24/5}\,t_Y)$ it is
$\bigl(\begin{smallmatrix}2g_2^2&-\sqrt{12/5}\,g_2g_Y\\-\sqrt{12/5}\,g_2g_Y&\frac65g_Y^2\end{smallmatrix}\bigr)$, with eigenvalues $0$ and $2g_2^2+\tfrac65g_Y^2$. p3-9 computes
the full $12\times12$ matrix on $\R^{512}$.
\end{proof}

So at the third rung the gauge-field candidate for the photon of the aligned vacua is the $\uu1_{\rm em}$ component of the internal connection, and the normal
$U(1)$ is a separate abelian field. What is established is a charge identification: a compact generator with the Standard-Model charges, anomaly free, which
commutes with the vacuum and couples to charged scalar fluctuations with the electron's unit. What is adopted is the Maxwell term (positive, since the
invariant form is positive on compact generators). Gauge redundancy, a conserved current, two physical polarisations, massless spin-one states and
quantisation then follow as in ordinary Yang--Mills theory for the adopted action: inherited, not derived from the Clifford data. That the field is the
physical photon also requires the electroweak stage, which \cref{p3:prop:neutral} shows the present potential does not supply (PH-06). On
the block that field couples to fermion number, and with $\su2_L$ gauged it cannot be gauged on that matter (\cref{p3:cor:notem}). The construction's
identification of its photon with the normal connection is therefore not compatible with the matter at $n=3$ (obstructed). What the second rung established survives
with $B_{\rm em}$ in place of $Q_A$: the massless criterion of \cref{cor:massless} at a Lorentz-invariant vacuum, the charged-scalar current, and a curved extension in which
the vacuum is parallel (\cref{p1:thm:curved}). The theorem rests on the choice of the gauged generator, and $B_{\rm em}$ is fixed only once the Standard-Model embedding of
\cref{sec:pstar} is adopted (constructed). The coupling of the second abelian field and the values of $f$ and of the gauge couplings are not decided (open).
\Cref{p3:thm:photon}(d) shows that the projector can itself carry an electroweak stage: its neutral doublet directions break $\mathfrak h_{\rm SM}$ to
$\su3_C\oplus\uu1_{\rm em}$ with the tree-level relation of one doublet \cite{Glashow61,Weinberg67}. Whether the vacuum moves there is decided by the potential, which
\cref{p3:sec:robust} examines.

% --- source: photon_selection/p3_identification#5 (The residual choice)
\subsection{The residual choice}
\label{p3:sec:residual}

By \cref{p1:prop:K,p1:thm:fixed}, with $\mathfrak h_{\rm SM}$ or $\mathfrak h_{\rm PS}$ gauged with positive couplings, the alignment energy vanishes on $F_3$ exactly at
$\Pi_A$, $\Pi_{-A}$, $\Pi_W$ and $\Pi_{-W}$. This subsection determines what separates these four points.

\begin{lemma}[Overlaps are Plücker functions; fixed]
\label{p3:lem:plucker}
For every oriented four-plane $P\subset V_+$, $\mathrm{vol}_P=\sum_Ip_I(P)\,e_I$, and for every operator $M$ on $\R^{512}$
\[
\Tr(\Pi_PM)=\tfrac12\Tr M+\tfrac12\sum_Ip_I(P)\Tr(e_IM).
\]
The four-blade components $\Tr(e_IM)$ vanish for $M=P_*$, $P_{\rm lep}$ and $\tfrac12(\id\pm\Gamma)$. The nonzero ones are $\Tr(e_A\Pfam)=64$,
$\Tr(e_AP_L)=-\Tr(e_WP_L)=32$, their negatives for $P_R$, $\Tr(e_AQ_v)=-\Tr(e_WQ_v)=64$ and $\Tr(e_A\Pi_A)=256$ (p3-10).
\end{lemma}

\begin{proof}
The frame vectors are orthonormal, so their Clifford product equals their exterior product, $\sum_I\det F_I\,e_I$. The blades $e_I$ are symmetric involutions with
$\Tr(e_Ie_J)=512\,\delta_{IJ}$. For the components it suffices to list index sets. The eight products of $\mathcal X$, $\mathcal Z$, $\mathcal U$ have the index sets $\emptyset$, $I\cup A\cup C$,
$A\cup W$, $\{8,9\}\cup W\cup C$ and their symmetric differences, none of which is a four-subset of $V_+$. With $\Gamma$, whose index set is all eighteen generators, the
products include $\Gamma\mathcal X=\chi_W$ and $\Gamma\mathcal X\mathcal Z=-\chi_W\mathrm{vol}_+=-\mathrm{vol}_A$, so $P_L=\tfrac1{16}(\id+\mathcal X)(\id+\mathcal Z)(\id+\mathcal U)(\id-\Gamma)$ contains
$\tfrac1{16}(\mathrm{vol}_A-\chi_W)$ and no other four-blade of $V_+$. The products of the $s_k$ have index sets of sizes $12$, $4$ inside $C$, and $8$ equal to $A\cup W$.
Among the products of the $X_i$ only $X_1=\mathrm{vol}_A$ is a four-blade of $V_+$ (\cref{p3:prop:restated}). $\gamma^0$ commutes with every $e_I$, which doubles the components
of $P_L$ in $Q_v$.
\end{proof}

\begin{proposition}[The overlaps; fixed]
\label{p3:prop:overlaps}
\begin{enumerate}[label=(\roman*),leftmargin=2.2em]
\item On all of $F_3$, $\Tr(\Pi P_*)=\Tr(\Pi P_{\rm lep})=32$ and $\Tr(\Pi\tfrac12(\id\pm\Gamma))=128$ are constant, while
\[
\Tr(\Pi_P\Pfam)=32(1+p_A)=\tfrac14\Tr(\Pi_P\Pi_A),\qquad \Tr(\Pi_PP_L)=16(1+p_A-p_W)=\tfrac12\Tr(\Pi_PQ_v),
\]
and $\Tr(\Pi_PP_R)=16(1-p_A+p_W)$ (p3-12).
\item $\Tr(\Pi P_*\Pi P_*)=\|P_*\Pi P_*\|^2\le\Tr(\Pi P_*)=32$, with equality iff $[\Pi,P_*]=0$. The points of $F_3$ commuting with $P_*$ are $\Pi_{\pm A}$ and $\Pi_{\pm W}$,
which are isolated, and the two components $\{\Pi_{U_A\oplus W'}\}$ and $\{\Pi_{A_0\oplus W'}\}$, $W'\in\widetilde{\mathrm{Gr}}_2(W)$, of dimension $4$ (p3-12, p3-13).
\item At the four points the linear and quadratic overlaps are those of \cref{p3:tab:overlaps} (p3-11).
\end{enumerate}
\end{proposition}

\begin{table}[t]
\centering
\small
\begin{tabular}{@{}lccccl@{}}
\toprule
$M$ & $\Pi_A$ & $\Pi_{-A}$ & $\Pi_W$ & $\Pi_{-W}$ & $\Tr(\Pi_PM)$ on $F_3$\\
\midrule
$P_*$ & $(32,32)$ & $(32,32)$ & $(32,32)$ & $(32,32)$ & $32$\\
$P_{\rm lep}$ & $(32,32)$ & $(32,32)$ & $(32,32)$ & $(32,32)$ & $32$\\
$\tfrac12(\id\pm\Gamma)$ & $(128,128)$ & $(128,128)$ & $(128,128)$ & $(128,128)$ & $128$\\
$P_L$ & $(32,32)$ & $(0,0)$ & $(0,0)$ & $(32,32)$ & $16(1+p_A-p_W)$\\
$P_R$ & $(0,0)$ & $(32,32)$ & $(32,32)$ & $(0,0)$ & $16(1-p_A+p_W)$\\
$Q_v$ & $(64,64)$ & $(0,0)$ & $(0,0)$ & $(64,64)$ & $32(1+p_A-p_W)$\\
$\Pfam$ & $(64,64)$ & $(0,0)$ & $(32,16)$ & $(32,16)$ & $32(1+p_A)$\\
\bottomrule
\end{tabular}
\caption{Overlaps $(\Tr\Pi M,\ \Tr\Pi M\Pi M)$ at the four aligned points, and the linear overlap on all of $F_3$ (p3-11, p3-12).}
\label{p3:tab:overlaps}
\end{table}

\begin{proof}
(i) is \cref{p3:lem:plucker} with $\Tr M$ and the components listed there. (ii) $P_*$ is a contraction, so $\|P_*\Pi P_*\|^2\le\|\Pi P_*\|^2=\Tr(P_*\Pi P_*)=32$, with equality
iff $(\id-P_*)\Pi P_*=0$, that is iff $\Pi P_*=P_*\Pi P_*$ is symmetric, that is iff $\Pi$ commutes with $P_*$. Every $\mathrm{vol}_P$ commutes with $J$, $\Gamma$ and
$\mathrm{vol}_+=-\mathcal Z$, so $\Pi_P$ commutes with $P_*$ iff $\mathrm{vol}_P$ commutes with $\chi_W$ and with $q_A$. $\Ad(\chi_W)$ acts on $V_+$ as the reflection $-1$ on $W$, $+1$ on $A$.
It fixes $\mathrm{vol}_P$ iff $P=(P\cap A)\oplus(P\cap W)$ with $\dim(P\cap W)$ even. With \cref{p1:thm:fixed} for $q_A$ this leaves $P=\pm A$, $P=\pm W$, and $P=A'\oplus W'$ with
$A'\in\{U_A,A_0\}$ and $W'$ a two-plane in $W$. p3-13 checks that the four points are isolated and that the components have dimension $4$. (iii) is the computation
p3-11.
\end{proof}

\begin{lemma}[A comass inequality; fixed]
\label{p3:lem:comass}
For every oriented four-plane $P\subset V_+=A\oplus W$, $|p_A(P)|+|p_W(P)|\le1$, with equality iff $P\in\{\pm A,\pm W\}$.
\end{lemma}

\begin{proof}
Let $M_A=(\langle a_i,f_j\rangle)$ and $M_W=(\langle w_i,f_j\rangle)$ for an oriented orthonormal frame $(f_j)$ of $P$, so that $M_A^{\!\top}M_A+M_W^{\!\top}M_W=\id$. The two
symmetric matrices commute, so in a common eigenbasis they are $\operatorname{diag}(\cos^2\theta_i)$ and $\operatorname{diag}(\sin^2\theta_i)$, with
$\theta_i\in[0,\tfrac\pi2]$ the principal angles between $P$ and $A$. Hence $|p_A|=\prod_i\cos\theta_i$ and $|p_W|=\prod_i\sin\theta_i$. For any two indices $i\neq j$,
\[
|p_A|+|p_W|\le\cos\theta_i\cos\theta_j+\sin\theta_i\sin\theta_j=\cos(\theta_i-\theta_j)\le1 .
\]
Equality for every pair forces all $\theta_i$ to equal one $\theta$, and then $\cos^4\theta+\sin^4\theta=1-2\cos^2\theta\sin^2\theta=1$ forces $\theta\in\{0,\tfrac\pi2\}$,
that is $P=\pm A$ or $P=\pm W$. p3-14 checks the identities at $300$ random planes and finds the maximum $1$ of $p_A-p_W$ only at $A$ and $-W$.
\end{proof}

\begin{corollary}[The extremes of the separating overlaps; fixed]
\label{p3:cor:extremes}
On $F_3$, $0\le\Tr(\Pi P_L)\le32$, with the maximum attained exactly at $\Pi_A$ and $\Pi_{-W}$ and the minimum exactly at $\Pi_{-A}$ and $\Pi_W$. Likewise
$0\le\Tr(\Pi\Pfam)\le64$, with the maximum only at $\Pi_A$ and the minimum only at $\Pi_{-A}$.
\end{corollary}

\begin{proof}
By \cref{p3:prop:overlaps}(i) and \cref{p3:lem:comass}, $|p_A-p_W|\le1$ with $p_A-p_W=1$ iff $P\in\{A,-W\}$ and $=-1$ iff $P\in\{-A,W\}$. Also $|p_A|\le1$ with equality
iff $P=\pm A$, since $|p_A|=\prod\cos\theta_i$.
\end{proof}

\begin{proposition}[A Klein group of the projector sector; fixed]
\label{p3:prop:klein}
Let $\mathsf C(\Pi)=\id-\Pi$, let $\mathsf D(\Pi)=\tfrac12(\id+\mathrm{vol}_+(2\Pi-\id))$ be the duality of \cref{p1:prop:duality}, and let $\mathsf E=\mathsf C\mathsf D$, so
that $\mathsf E(\Pi)=\tfrac12(\id+\mathrm{vol}_+)-\mathrm{vol}_+\Pi$.
\begin{enumerate}[label=(\roman*),leftmargin=2.2em]
\item $\{1,\mathsf C,\mathsf D,\mathsf E\}$ is a Klein four-group of $K$-equivariant involutions of $F_3$, with $\mathsf E(\Pi_P)=\Pi_{-P^\perp}$. It acts simply transitively on
the four points: $\mathsf D$ exchanges $\Pi_A$ and $\Pi_W$, $\mathsf C$ exchanges $\Pi_A$ and $\Pi_{-A}$, and $\mathsf E$ exchanges $\Pi_A$ and $\Pi_{-W}$ (p3-16).
\item Each element preserves every mass matrix $\Tr([T,\Pi][T',\Pi])$, $T,T'\in\mathfrak k$, and, for connections valued in $\mathfrak k$ and the interior Lorentz algebra, the
kinetic term and the current $j_\Pi$ of every $Q\in\mathfrak k$ (p3-16).
\item $\mathsf E$ is the identity on the half $\mathcal Z=+1$, the range of $\tfrac12(\id+\mathcal Z)$, which contains $P_*$: $\mathsf E(\Pi)\tfrac12(\id+\mathcal Z)=\Pi\tfrac12(\id+\mathcal Z)$, and
$P_*\mathsf E(\Pi)P_*=P_*\Pi P_*$. On that half $\Pi_A=\Pi_{-W}=\tfrac12(\id-\chi_W)$ and $\Pi_W=\Pi_{-A}=\tfrac12(\id+\chi_W)$, and there $t^L_i=\Pi_A\tau_i$ and
$t^R_i=\Pi_W\tau_i$. On the half $\mathcal Z=-1$, $\Pi_A=\Pi_W$ and $\Pi_{-A}=\Pi_{-W}$ (p3-17).
\end{enumerate}
Hence no alignment energy, kinetic term or current separates any two of the four points. No coupling that depends on $\Pi$ only through $\Pi\tfrac12(\id+\mathcal Z)$ separates
$\Pi_A$ from $\Pi_{-W}$, or $\Pi_{-A}$ from $\Pi_W$.
\end{proposition}

\begin{proof}
$\mathrm{vol}_+$ is $K$-invariant, commutes with the interior gammas and with every $\Pi_P$, and squares to $\id$ (\cref{p1:lem:census,p1:prop:duality}). $\mathsf C$ and
$\mathsf D$ commute, since both $\mathsf C\mathsf D(\Pi)$ and $\mathsf D\mathsf C(\Pi)$ equal $\tfrac12(\id-\mathrm{vol}_+(2\Pi-\id))$. (ii) For $T$ in $\mathfrak k$ or in the
interior Lorentz algebra, $[T,\mathsf C\Pi]=-[T,\Pi]$, $[T,\mathsf D\Pi]=\mathrm{vol}_+[T,\Pi]$ and $[T,\mathsf E\Pi]=-\mathrm{vol}_+[T,\Pi]$. Since $\mathrm{vol}_+$ is orthogonal and
commutes with $[T,\Pi]$, the traces are preserved. $\mathrm{vol}_+$ is parallel for such connections, so $D_\mu$ transforms in the same way. (iii) On $\mathcal Z=+1$,
$\mathrm{vol}_+=-\id$, so $\mathsf E(\Pi)=\Pi$ there. $\mathrm{vol}_A=\mathrm{vol}_+\chi_W$ equals $-\chi_W$ there and $+\chi_W$ on $\mathcal Z=-1$; with
$t^L_i=\tfrac12(\id-\chi_W)\tau_i$ \cref{thm:PSalg} the remaining identities follow.
\end{proof}

On the half of the module that contains the matter, the vacuum $\Pi_A$ is the support projector of $\su2_L$, and $\Pi_W$ that of $\su2_R$. The residual choice is
therefore, on the matter, the choice of which weak $\su2$ acts inside the image of the vacuum. It is not a choice of spacetime chirality: the chirality halves have
constant overlap, and the vector-like $Q_v$ separates exactly as the chiral $P_L$ does (\cref{p3:tab:overlaps}).

\begin{proposition}[Normalisers; fixed]
\label{p3:prop:normaliser}
\begin{enumerate}[label=(\roman*),leftmargin=2.2em]
\item The fixed vectors of $\rho_+(\mathfrak h_{\rm SM})$ and of $\rho_+(\mathfrak h_{\rm PS})$ in $V_+$ are exactly $A$. A reflection of $W$ conjugates $\su2_L$ onto $\su2_R$.
Hence every element of $K$ normalising $\mathfrak h_{\rm SM}$ fixes each of $\Pi_A$, $\Pi_{-A}$, $\Pi_W$, $\Pi_{-W}$ (p3-18).
\item $g=E_9E_{15}=\exp(\tfrac\pi2E_9E_{15})$ normalises $\mathfrak h_{\rm PS}$, exchanges $\su2_L$ and $\su2_R$, and sends $\mathrm{vol}_A\mapsto-\mathrm{vol}_A$,
$\mathrm{vol}_W\mapsto-\mathrm{vol}_W$ and $\mathcal X\mapsto-\mathcal X$. So the normaliser of $\mathfrak h_{\rm PS}$ identifies $\Pi_A$ with $\Pi_{-A}$ and $\Pi_W$ with $\Pi_{-W}$, by
elements that do not preserve $P_*$. No element of either normaliser maps $\Pi_{\pm A}$ to $\Pi_{\pm W}$ (p3-18).
\item $g_{AW}=\prod_i\tfrac1{\sqrt2}(\id+a_iw_i)$, with $(a_i)=(E_2,E_3,E_8,E_9)$ and $(w_i)=(E_7,E_{13},E_{14},E_{15})$, is an element of $\mathrm{Spin}(8)$ with
$\Ad(g_{AW})\mathrm{vol}_A=\mathrm{vol}_W$ and $\Ad(g_{AW})\mathrm{vol}_W=\mathrm{vol}_A$, which normalises neither $\mathfrak h_{\rm SM}$ nor $\mathfrak h_{\rm PS}$. Moreover
$\mathrm{vol}_A\mathrm{vol}_W=\mathrm{vol}_+$ (p3-18).
\end{enumerate}
\end{proposition}

\begin{proof}
(i) $\Ad(k)$ acts on $F_3$ through the rotation $R=\rho(k)$ of $V_+$ (\cref{p1:thm:classify}). If $\Ad(k)$ normalises $\mathfrak h$, $R$ preserves the fixed space $A$ of
$\rho_+(\mathfrak h)$ and its complement $W$. For $\mathfrak h_{\rm SM}$, $R|_W$ normalises the derived algebra $\su2_L$ of $\rho_+(\mathfrak h_{\rm SM})$. A rotation of $W$ of
determinant $-1$ is a rotation times a reflection, and it maps $\su2_L$ onto $\su2_R$; so $\det R|_W=1$, hence $\det R|_A=1$, and $R$ fixes the oriented planes $A$ and $W$. (ii)
$\Ad(g)$ is the rotation by $\pi$ in the plane $(E_9,E_{15})$, of determinant $-1$ on $A$ and on $W$; it preserves $W$ and $V_-$, hence $\mathfrak h_{\rm PS}$. $\Gamma$ commutes
with $g$ and $\Ad(g)\chi_W=-\chi_W$. The last statement follows from (i), which applies to $\mathfrak h_{\rm PS}$ as well. (iii) $\Ad(\exp\tfrac\pi4a_iw_i)$ rotates the plane
$(a_i,w_i)$ by $\tfrac\pi2$; the signs and the failure to normalise are computed in p3-18.
\end{proof}

With $\mathfrak h_{\rm SM}$ gauged the four points are therefore distinct points of the gauged theory, although the projector sector cannot tell them apart. With
$\mathfrak h_{\rm PS}$ gauged, $\Pi_A$ and $\Pi_{-A}$ are related by a left--right exchange of the Pati--Salam algebra, which the matter block breaks.

% --- source: photon_selection/p3_identification#6 (Robustness)
\subsection{Robustness}
\label{p3:sec:robust}

\begin{lemma}[Invariant couplings keep fixed critical points; inherited, with proof]
\label{p3:lem:robust}
Let $H\subset K$ be connected with Lie algebra $\mathfrak h$, let $f$ be a smooth $H$-invariant function on $F_3$, and let $\Pi_0\in F_3$ be fixed by $H$. Then $df_{\Pi_0}$ is
an $H$-invariant covector. If $T_{\Pi_0}F_3$ contains no nonzero $\mathfrak h$-invariant vector, $\Pi_0$ is a critical point of $f$. In particular $\Pi_0$ stays critical
for $V_0+tf$ for every $t$ whenever $V_0$ is $H$-invariant too, and it stays a nondegenerate local minimum as long as the Hessian of $V_0+tf$, an $H$-invariant
symmetric form, stays positive definite.
\end{lemma}

\begin{proof}
For $h\in H$, $f\circ h=f$ and $h\Pi_0=\Pi_0$ give $df_{\Pi_0}\circ dh_{\Pi_0}=df_{\Pi_0}$. The metric $\Tr(d\Pi\,d\Pi)$ is $K$-invariant (\cref{p1:lem:vector}), so the gradient
is an $H$-invariant tangent vector, and for a connected group the invariant vectors are those killed by $\ad\mathfrak h$ \cite{Michel71,Palais79}. The Hessian at a critical
point is invariant for the same reason.
\end{proof}

\begin{proposition}[Invariant tangent vectors; fixed]
\label{p3:prop:invariants}
The dimensions of the $\mathfrak h$-invariant subspaces of the tangent spaces at $\Pi_A$ and $\Pi_W$, and so at $\Pi_{-A}$ and $\Pi_{-W}$, whose tangent spaces are the same
up to sign, are: $\mathfrak h_{\rm SM}$, $\mathfrak h_{\rm PS}$, $\su2_L$ and $\uu1_Y\oplus\su3_C$: $0$; $\uu1_{\rm em}\oplus\su3_C$: $8$, the neutral components $A\otimes W_0$;
$\uu1_A$: $8$; $\uu1_A\oplus\uu1_{\rm em}$: $4$ (p3-19).
\end{proposition}

\begin{corollary}[Matter couplings that respect $\mathfrak h$; fixed, with an inherited principle]
\label{p3:cor:matter}
Let the vacuum energy on $F_3$ be the sum of an alignment energy and of couplings to matter that respect $\mathfrak h\in\{\mathfrak h_{\rm SM},\mathfrak h_{\rm PS}\}$: functions of $\Pi$
and of operators fixed by $\mathfrak h$, such as the overlaps of \cref{p3:prop:overlaps}, Yukawa-type kernels commuting with $\mathfrak h$, or the fermion loop of a Dirac operator
with $\mathfrak h$-invariant couplings \cref{sec:matter}. Such an energy is $H$-invariant, since $f(\Ad(h)\Pi;M)=f(\Pi;\Ad(h^{-1})M)$, and $\Pi_{\pm A}$, $\Pi_{\pm W}$ remain
critical points of it. All linear overlaps of \cref{p3:tab:overlaps}, including the one with $\Pfam$, which does not commute with $\su3_C$, are $\mathfrak h_{\rm SM}$-invariant
functions on $F_3$, and their first variations vanish at the four points along all of $\mathfrak k$ (p3-20). With only $\uu1_{\rm em}\oplus\su3_C$ gauged the eight neutral
directions are invariant, and an invariant coupling can move the vacuum along the $\uu1_{\rm em}$-fixed component $\widetilde{\mathrm{Gr}}_4(A\oplus W_0)$ of
\cref{p1:thm:fixed}.
\end{corollary}

\begin{proof}
\Cref{p3:lem:robust} with \cref{p3:prop:invariants}. The invariance of the overlaps is p3-20; for $\Pfam$ it also follows from $\Tr(\Pi_P\Pfam)=\tfrac14\Tr(\Pi_P\Pi_A)$ and
$[\mathfrak h_{\rm SM},\Pi_A]=0$.
\end{proof}

What invariant couplings can do is change the Hessian. Write $\Pi(x)=\Ad(e^X)\Pi_0$ with $X=\sum x_{aw}\tfrac12aw$, as in \cref{p1:thm:tangent}, and
\[
c_{\rm SM}=6g_2^2+\tfrac65g_Y^2,\qquad c_{\rm PS}=12g^2 .
\]

\begin{proposition}[Hessians at the four points; fixed]
\label{p3:prop:hessian}
At each of the four points $V_{\mathfrak h_{\rm SM}}(\Pi(x))=128\,c_{\rm SM}|x|^2+O(x^3)$ and $V_{\mathfrak h_{\rm PS}}(\Pi(x))=128\,c_{\rm PS}|x|^2+O(x^3)$. The overlaps expand as
\[
\Tr(\Pi(x)P_L)=\begin{cases}32-8|x|^2+O(x^3)&\text{at }\Pi_A,\ \Pi_{-W},\\ 8|x|^2+O(x^3)&\text{at }\Pi_{-A},\ \Pi_W,\end{cases}
\]
\[
\Tr(\Pi(x)\Pfam)=\begin{cases}64-16|x|^2+\cdots&\text{at }\Pi_A,\\ 16|x|^2+\cdots&\text{at }\Pi_{-A},\\ 32+O(x^4)&\text{at }\Pi_{\pm W}\end{cases}
\]
(p3-21). For $V=\xi_1 V_{\mathfrak h_{\rm SM}}-\xi_2\Tr(\Pi P_L)-\xi_3\Tr(\Pi\Pfam)$, $\xi_1,\xi_2,\xi_3\ge0$, every Hessian at the four points is a multiple of the
identity, with coefficient $128\xi_1 c_{\rm SM}+h$, where $h=8\xi_2+16\xi_3$ at $\Pi_A$, $8\xi_2$ at $\Pi_{-W}$, $-8\xi_2$ at $\Pi_W$ and $-8\xi_2-16\xi_3$ at $\Pi_{-A}$.
Hence $\Pi_A$ is a nondegenerate local minimum unless $\xi_1=\xi_2=\xi_3=0$, and $\Pi_{-W}$ is one iff $\xi_1>0$ or $\xi_2>0$. At $\xi_1=\xi_2=0<\xi_3$ the
Hessian at $\Pi_{-W}$ vanishes, and $\Pi_{-W}$ is \emph{not} a local minimum: along the curve that turns each $w_j$ towards $a_j$ through the same angle $t$, with
orientation such that $p_A=\sin^4t$, $V=-32\xi_3(1+\sin^4t)<V(\Pi_{-W})$ (z2-3). $\Pi_W$ is a nondegenerate local minimum iff $\xi_2<16\xi_1 c_{\rm SM}$, and $\Pi_{-A}$ iff
$\xi_2+2\xi_3<16\xi_1 c_{\rm SM}$; for the opposite strict inequalities all sixteen directions are unstable at once. On the boundaries the Hessian vanishes, and the
quadratic test decides nothing (PH-02).
\end{proposition}

\begin{proof}
Since $[t_a,\Pi_0]=0$, $\|[t_a,\Pi(x)]\|^2=\|[[t_a,X],\Pi_0]\|^2+O(x^3)$. $[t_a,X]$ corresponds to $x\,\rho_+(t_a)^{\!\top}$ on $A\otimes W$, so the sum is
$128\,\tr(x\,C\,x^{\!\top})$ with $C=\sum_ag_a^2\rho_+(t_a)\rho_+(t_a)^{\!\top}$ on $W$. For $\mathfrak h_{\rm SM}$ the Casimir of $\su2_L$ and $\rho_+(t_Y)^2=-\tfrac14$ give
$C=(6g_2^2+\tfrac65g_Y^2)\id_W$, and $\su3_C$ gives nothing; for $\mathfrak h_{\rm PS}$ each of the four directions of $W$ lies in three unit blades, so $C=12g^2\id_W$. The
overlaps follow from \cref{p3:prop:overlaps}(i) with $p_A(\Pi(x))=1-\tfrac12|x|^2+O(x^4)$ near $A$, and $p_W=\det(O(x))=O(x^4)$ there; near $\pm W$ the roles are exchanged.
p3-21 checks the second derivatives $\Tr([X,[X,\Pi_0]]P)$ exactly and by finite differences. The stability conditions compare $128\xi_1 c_{\rm SM}$ with $8\xi_2$, and with
$8\xi_2+16\xi_3$ at $\Pi_{-A}$.
\end{proof}

The sign of $\xi_2$ is a choice. For $\xi_2<0$ the roles of $\{\Pi_A,\Pi_{-W}\}$ and $\{\Pi_{-A},\Pi_W\}$ are exchanged. The Hessians found here are the multiples of
the identity in the sixteen-dimensional space of $\mathfrak h_{\rm SM}$-invariant forms of \cref{p1:cor:critical}; a coupling with a component of the form
$N\otimes\omega_Y$ would split the doublet directions.

% --- source: photon_selection/p3_identification#7 (The complete selection)
\subsection{The complete selection}
\label{p3:sec:selection}

The alignment result used is \cref{p1:prop:K,p1:thm:fixed}: for positive couplings $V_{\mathfrak h_{\rm SM}}\ge0$ on $F_3$, with zero set exactly $\{\Pi_{\pm A},\Pi_{\pm W}\}$.

\begin{theorem}[Selection on $F_3$]
\label{p3:thm:selection}
Assume:
\begin{enumerate}[label=(H\arabic*),leftmargin=2.6em]
\item the vacuum projector is a point of $F_3$ (constructed: \cref{p1:sec:answer});
\item $\mathfrak h_{\rm SM}$ is gauged, embedded as in \cref{sec:pstar}, with couplings $g_3,g_2,g_Y>0$ (constructed);
\item the sign hypothesis: the gauge fluctuations contribute $\xi_1 V_{\mathfrak h_{\rm SM}}$ to the vacuum energy with $\xi_1\ge0$ (inherited mechanism of vacuum alignment
\cite{Peskin80,Preskill81}; its sign is proved for vector-like theories \cite{Witten83} and assumed here);
\item the vacuum couples to the block by $-\xi_2\Tr(\Pi P_L)$ with $\xi_2>0$ (constructed, sign chosen);
\item optionally, it couples to the family block by $-\xi_3\Tr(\Pi\Pfam)$ with $\xi_3>0$ (constructed).
\end{enumerate}
Let $V=\xi_1 V_{\mathfrak h_{\rm SM}}-\xi_2\Tr(\Pi P_L)-\xi_3\Tr(\Pi\Pfam)$, with $\xi_3=0$ when (H5) is not assumed.
\begin{enumerate}[label=(\alph*),leftmargin=2.2em]
\item Under (H1)--(H3) with $\xi_1>0$ the minimum set is $\{\Pi_{\pm A},\Pi_{\pm W}\}$. At $\xi_1=0$, $V\equiv0$ and every point of $F_3$ is a minimum (PH-01). Overlaps with $P_*$, $P_{\rm lep}$ and $\tfrac12(\id\pm\Gamma)$, linear or quadratic, take equal values at
these points and do not separate them.
\item Under (H1)--(H4) the global minimum value is $-32\xi_2$, attained exactly at $\Pi_A$ and $\Pi_{-W}$, for every $\xi_1\ge0$. These two points coincide on the half
$\mathcal Z=+1$ that contains $P_*$, and they are exchanged by $\mathsf E$, which preserves every term of $V$ except the one of (H5). $\Pi_W$ and $\Pi_{-A}$ remain critical, and are
nondegenerate local minima iff $\xi_2<16\xi_1 c_{\rm SM}$ (with $\xi_3=0$; boundary cases as in \cref{p3:prop:hessian}).
\item Under (H1)--(H5) the global minimum is $\Pi_A$ alone, with value $-32\xi_2-64\xi_3$.
\item At the selected vacuum all of $\mathfrak h_{\rm SM}$ is unbroken. $\Pi_A$ is the electroweak-symmetric point, not the observed electroweak-broken vacuum:
the generator $B_{\rm em}$ is massless there, but so are the $W^\pm$ and the $Z$, and a photon distinguished from them by mass needs the electroweak stage of
\cref{p3:thm:photon}(d), which the potential $V$ disfavours (\cref{p3:prop:neutral}). The field along $Q_A$ is a second abelian field (PH-03).
\end{enumerate}
\end{theorem}

\begin{proof}
(a) is \cref{p1:prop:K,p1:thm:fixed} and \cref{p3:tab:overlaps}. (b) $V\ge\xi_1\cdot0-32\xi_2$ by \cref{p3:cor:extremes}, with equality iff $V_{\mathfrak h_{\rm SM}}=0$ and
$\Tr(\Pi P_L)=32$, that is iff $\Pi\in\{\Pi_A,\Pi_{-W}\}$, where both conditions hold. The coincidence and the role of $\mathsf E$ are \cref{p3:prop:klein}; the overlap with $P_L$
depends on $\Pi$ only through $\Pi\tfrac12(\id+\mathcal Z)$. Stability is \cref{p3:prop:hessian}. (c) With $\xi_3>0$ equality requires in addition $\Tr(\Pi\Pfam)=64$, which holds only at
$\Pi_A$ (\cref{p3:cor:extremes}). (d) is \cref{p1:thm:vacuum} and \cref{p3:thm:photon}. p3-22 minimises $V$ from random starts for several couplings, including
$\xi_2>16\xi_1 c_{\rm SM}$ and $\xi_1=0$, and finds exactly these minima.
\end{proof}

Two features of the theorem could have come out otherwise. The overlap with $P_*$ itself, which the matter block suggests first, is constant on $F_3$ and selects nothing,
and the selection is made by the weak-chirality half $P_L$. Once that overlap is present the alignment energy is not needed for the selection; it only stabilises the
other two critical points. The pair $\{\Pi_A,\Pi_{-W}\}$ is one projector on the half of the module that carries the matter, so for physics on the block the selected
vacuum is unique (interpretive). A unique point of $F_3$ needs a coupling that sees the other half. The overlap with $\Pfam$ is such a coupling, but $\Pfam$ carries no colour,
and adopting it amounts to choosing $\Pi_A$, since $\Tr(\Pi\Pfam)=\tfrac14\Tr(\Pi\Pi_A)$ on $F_3$ (constructed). The rank-$256$ projector $\Pi_A$ is the projector $P_1$ of the
projector family \cref{thm:table}, on which colour and a chiral doublet first coexist \cref{thm:su3,thm:p1chiral}.

\begin{proposition}[The potential opposes the electroweak deformation; fixed]
\label{p3:prop:neutral}
Let $\Pi_\theta=\Ad(\exp\tfrac\theta2E_8E_{14})\Pi_A$, the neutral deformation of \cref{p3:thm:photon}(d). Then $\Tr(\Pi_\theta P_L)=16(1+\cos\theta)$,
$\Tr(\Pi_\theta\Pfam)=32(1+\cos\theta)$, $V_{\mathfrak h_{\rm SM}}(\Pi_\theta)>0$ for $\theta\notin\pi\Z$, and
\[
V(\Pi_\theta)-V(\Pi_A)=\xi_1 V_{\mathfrak h_{\rm SM}}(\Pi_\theta)+(16\xi_2+32\xi_3)(1-\cos\theta)\ \ge0 .
\]
Under the positive coefficients of \cref{p3:thm:selection} every term favours the aligned, electroweak-symmetric point (z2-5).
\end{proposition}

\begin{proof}
$\Ad(\exp\tfrac\theta2E_8E_{14})$ turns $E_8$ into $\cos\theta\,E_8+\sin\theta\,E_{14}$, so $p_A(\Pi_\theta)=\cos\theta$ and $p_W(\Pi_\theta)=0$; apply \cref{p3:prop:overlaps}(i).
The mass matrix of \cref{p3:thm:photon}(d) is nonzero for $\theta\notin\pi\Z$.
\end{proof}

So the electroweak stage that \cref{p3:thm:photon}(d) exhibits is not reached by this potential. A transition to $\Pi_\theta$ needs a further contribution with the
opposite sign along the neutral doublet direction, for instance a negative mass term for the doublets. No such term is derived here, and none is adopted:
the mechanism of electroweak breaking is missing (PH-04).

\begin{remark}[Status of the selecting potential and of the scalar fluctuations]
\label{p3:rem:status}
(a) \emph{Origin} (PH-07). The alignment term rests on the sign hypothesis (H3). The overlaps with $P_L$ and $\Pfam$, and their signs, are adopted. Since
$\Tr(\Pi\Pfam)=\tfrac14\Tr(\Pi\Pi_A)$ on $F_3$, the $\xi_3$-term encodes its target. \Cref{p3:thm:selection} is therefore a conditional selection given an
adopted potential, not a first-principles selection mechanism. (b) \emph{The four doublets} (PH-08). At each aligned point the sixteen real tangent directions
form four complex doublets of $\su2_L$ (\cref{p1:thm:tangent}). That is representation content. If all sixteen modes propagated and an electroweak stage removed
three as longitudinal gauge modes, thirteen real scalars would remain. Their masses (with $V$ above, all sixteen share one Hessian coefficient at $\Pi_A$), their
mixing, their Yukawa couplings and any constraint that removes them are not determined, so no scalar phenomenology is claimed.
\end{remark}

\begin{proposition}[Restated data of the projectors; restated]
\label{p3:prop:restated}
$X_1=J^{(1)}J^{(2)}=\mathrm{vol}_A$, so $P_1=\Pi_A$, with $\dim\mathfrak g_{\rm phys}(P_1)=51$, $\dim\mathfrak k(P_1)=27$ and a three-dimensional annihilator, as in
\cref{thm:table}, and $\Pfam\le\Pi_A$. The projector $P_{\rm lep}$ of \cref{p3:def:plep} has rank $64$, commutes with the interior gammas and with $J$, and with $P_*$,
$\Pfam$ and $Q_v$ reproduces the alignment energies $3$, $5$, $9$, $5$ and stabilisers $37$, $23$, $9$, $25$ of \cref{prop:alignment} (p3-15).
\end{proposition}

\begin{proof}
$X_1=e^{(1)}_2e^{(1)}_3e^{(2)}_2e^{(2)}_3=E_2E_3E_8E_9$. The other statements are the computation p3-15.
\end{proof}

% --- source: photon_selection/p3_identification#8 (The answer)
\subsection{The answer}
\label{p3:sec:answer}

\begin{enumerate}[label=(\alph*),leftmargin=2.2em]
\item \emph{Which $\uu1$ is the photon.} At the third rung the photon of the aligned vacua is the $\uu1_{\rm em}$ gauge field of the internal connection, generated by
$B_{\rm em}=-\tfrac12E_7E_{13}+\tfrac16H_C$. It is massless at $\Pi_{\pm A}$ and $\Pi_{\pm W}$, its projector current is carried by four complex scalars of the electron's
charge magnitude, and on the block it couples with the Standard-Model charges without anomaly (fixed, given the adopted embedding, which is constructed).
\item \emph{The construction's charge.} $Q=\tfrac12J_1$ becomes $Q_A=\tfrac12E_2E_3$, whose charge operator is $\tfrac12\mathcal U$. On the block it is half the fermion number, the
same on every state, and anomalous with $\su2_L$ (fixed; anomaly restated). It is not electromagnetism (obstructed). The normal $SO(2)$ connection that carries it is a
separate abelian field (fixed), and its coupling is not decided (open).
\item \emph{The counts and the dichotomy.} Of the three counts against the normal $U(1)$, the spinless count is unchanged, the adapted-generator count is removed for both
charges, and the sign count changes from forced to adopted; a fourth count, the charge on the block, decides. Both alternatives of the dichotomy of
\cref{sec:transport} hold at $n=3$ (fixed; the Kaluza--Klein reading of the sign count is interpretive).
\item \emph{Which vacuum.} Alignment by $\mathfrak h_{\rm SM}$ leaves $\Pi_{\pm A},\Pi_{\pm W}$ (fixed, \cref{p1:thm:fixed}). A Klein group of symmetries of the projector sector
permutes them simply transitively; no alignment energy, kinetic term or current separates them (fixed). The weak-chirality overlap $-\xi_2\Tr(\Pi P_L)$, $\xi_2>0$, selects
$\{\Pi_A,\Pi_{-W}\}$, which is one projector on the matter half (fixed under constructed hypotheses). A single point requires a coupling off that half, such as
$-\xi_3\Tr(\Pi\Pfam)$, which selects $\Pi_A$ (constructed). With only alignment ($\xi_2=\xi_3=0$) the four points need $\xi_1>0$, and at $\xi_1=0$ every point
is a minimum (PH-01). The selected $\Pi_A$ preserves all of $\mathfrak h_{\rm SM}$: it is the electroweak-symmetric point, not the observed vacuum (PH-03), and the
adopted potential opposes the neutral electroweak deformation (\cref{p3:prop:neutral}; PH-04).
\item \emph{Robustness.} $\mathfrak h_{\rm SM}$- or $\mathfrak h_{\rm PS}$-invariant couplings, including matter couplings, keep the four points critical, because their
tangent spaces carry no invariant vector (fixed; principle inherited). They can only change the Hessians, which for the couplings considered are multiples of the identity;
$\Pi_W$ and $\Pi_{-A}$ lose stability for $\xi_2>16\xi_1 c_{\rm SM}$ and $\xi_2+2\xi_3>16\xi_1 c_{\rm SM}$, and $\Pi_{-W}$ is not a local minimum at
$\xi_1=\xi_2=0<\xi_3$ (fixed; corrected, PH-02).
\item \emph{Open.} The dynamical origin and sign of the overlap couplings, the scales $f$, $g$, $\xi_1$, $\xi_2$, $\xi_3$, the electroweak stage and its mechanism, the coupling of the second abelian field and how its anomaly is
excluded, cancelled or decoupled (\cref{p3:rem:anomaly}), whether the
half $\mathcal Z=-1$ is physical, and the origin of the Maxwell term of $\uu1_{\rm em}$ are open.
\end{enumerate}

\begin{table}[t]
\centering
\footnotesize
\begin{tabular}{@{}>{\raggedright\arraybackslash}p{8.4cm}>{\raggedright\arraybackslash}p{4.3cm}>{\raggedright\arraybackslash}p{2.2cm}@{}}
\toprule
statement & class & checks\\
\midrule
\cref{p3:def:plep}, \cref{p3:prop:restated}: $P_{\rm lep}$, $P_1=\Pi_A$, the table of \cref{prop:alignment} & restated (\cref{thm:table}, \cref{prop:alignment}) & p3-15\\
\cref{p3:prop:charges}: $\hat Q_A=\tfrac12\mathcal U$, $q_AP_*=JP_*$, charge table, Weyl block & fixed; $q_AP_*=JP_*$ and the $\hat Y$, $\hat Q_{\rm em}$ columns restated & p3-1 to p3-3\\
\cref{p3:cor:notem}: $Q_A$ is $\tfrac12F$ on the block, not the electric charge, anomalous; $\uu1_{\rm em}$ anomaly free & fixed; obstructed; anomaly restated; anomaly conditions inherited & p3-4, p3-5, z2-1\\
\cref{p3:rem:anomaly}: the normal connection must be excluded, cancelled (needs the mirror matter) or decoupled & fixed (cancellation trace); options open & z2-2\\
\cref{p3:tab:counts}: the four counts and the dichotomy & restated (i); fixed (ii), (iv); constructed and interpretive (iii) & p3-4, p3-5\\
\cref{p3:prop:jointfixed}: joint fixed points of $Q_A$ and $B_{\rm em}$ & fixed & p3-6\\
\cref{p3:prop:torus}: torus weights of the fluctuations & fixed & p3-7\\
\cref{p3:thm:photon}: the photon candidate is $\uu1_{\rm em}$ (charge identification); common charge unit; $Q_A$ field separate; $m_W^2/m_Z^2$; $\sin^2\theta_W=\tfrac38$ by normalisation and $g_Y=g_2$ & fixed; gauged generator constructed; Maxwell term adopted; photon dynamics inherited, not derived; (b) obstructed & p3-4, p3-8, p3-9, z2-6\\
\cref{p3:lem:plucker}, \cref{p3:prop:overlaps}: overlaps as Plücker functions, values, commutant of $P_*$ & fixed & p3-10 to p3-13\\
\cref{p3:lem:comass}, \cref{p3:cor:extremes}: comass inequality, extremes of the separating overlaps & fixed & p3-12, p3-14\\
\cref{p3:prop:klein}: Klein group, $\mathsf E$ trivial on the matter half & fixed & p3-16, p3-17\\
\cref{p3:prop:normaliser}: normalisers of $\mathfrak h_{\rm SM}$, $\mathfrak h_{\rm PS}$; $g_{AW}$ & fixed & p3-18\\
\cref{p3:lem:robust}: invariant couplings keep fixed critical points & inherited \cite{Michel71,Palais79}, proof given & ---\\
\cref{p3:prop:invariants}, \cref{p3:cor:matter}: invariant tangent vectors, matter couplings & fixed; invariance of induced energies inherited & p3-19, p3-20\\
\cref{p3:prop:hessian}: Hessians and stability conditions; $\Pi_{-W}$ not a minimum at $\xi_1=\xi_2=0<\xi_3$; boundary cases degenerate & fixed (corrected, PH-02) & p3-21, z2-3, z2-4\\
\cref{p3:thm:selection}: selection of $\{\Pi_A,\Pi_{-W}\}$, and of $\Pi_A$ with (H5) & fixed under (H1), (H2), (H4), (H5) constructed and (H3) inherited; (a) needs $\xi_1>0$ (PH-01); conditional selection, not first-principles & p3-22\\
\cref{p3:prop:neutral}: the potential opposes the electroweak deformation & fixed; electroweak mechanism missing & z2-5\\
\cref{p3:rem:status}: origin of the potential adopted; four doublets as representation content only & constructed; open & ---\\
unique vacuum for physics on the block; Kaluza--Klein reading of the sign count & interpretive & ---\\
origin and sign of the overlap couplings; scales; the second abelian field & open & ---\\
\bottomrule
\end{tabular}
\caption{Claims of \cref{p3:sec:main} and their classes. The checks are the \texttt{PASS} lines printed by \texttt{p3\_identification.py}.}
\label{p3:tab:claims}
\end{table}

% --- source: tex/s26_photon#0 ((head))


% --- source: tex/s26_photon#1 (Sources, and whether the photon can be massless)
\section{Sources, and whether the photon can be massless}
\label{sec:photon}

\Cref{thm:twoops} fixes the form of Maxwell's equations and leaves the source entirely open. The geometry of the curved arena supplies two candidate
currents, and testing them gives a sharp negative result about the photon, together with the criterion that repairs it and a vacuum, at the second rung, that
satisfies the criterion.

% --- source: tex/s26_photon#2 (Two currents)
\subsection{Two currents}

On the slice $\Sigma$ the oriented normal two-plane bundle of \cref{sec:curvedmodel} has an $SO(2)$ connection $A$ with curvature $F_N=dA$, an ordinary two-form
on spacetime whose six components sit in the module of \cref{thm:emfield}, and which is distinct from the compressed spinor curvature $\mathcal F_{\rm obs}$.
Because the structure group is abelian, $dF_N=0$ identically: the homogeneous half of Maxwell's equations is the Bianchi identity of the normal bundle, not an
assumption. Setting $j^\nu_{\rm geom}:=\lambda\,\nabla_\mu F_N^{\mu\nu}$ gives a current conserved identically, by the antisymmetry of $F_N$. On a flat base the
connection $A=Bx\,dz$ has constant curvature and zero current, while $A=a(x)\,dz$ gives $F_N=a'(x)\,dx\wedge dz$ and $j^z_{\rm geom}=\lambda\,a''(x)$. This current is
defined from $F_N$, so it is the current of a geometry that has been specified, not an independent source; its value is to show that continuity can be met by a
geometric object.

The stronger candidate is independent of $F$. Fix a compact $\uu1$ generator $Q$, let $D_\mu\Pi=\nabla_\mu\Pi+A_\mu[Q,\Pi]$, and add a Maxwell term
$-F_{\mu\nu}F^{\mu\nu}/4g_A^2$. Varying the kinetic term $\tfrac{f^2}2\Tr(D_\mu\Pi D^\mu\Pi)$ with respect to $A_\mu$ gives
\begin{equation}
\label{eq:jPi}
j^\mu_\Pi=-f^2\Tr\bigl(Q\,[\Pi,D^\mu\Pi]\bigr)=-f^2\Tr\bigl([Q,\Pi]\,D^\mu\Pi\bigr),
\end{equation}
the component of projector motion along the chosen charge rotation. It is conserved on the projector equation of motion, $D_\mu\mathcal C^\mu=0$, rather than
identically.

\begin{theorem}[The paired-mode projector sources a massive vector]
\label{thm:photonmass}
For $Q=J_N/2$ and the angle ansatz of \cref{ex:stueckelberg}, $\Tr([Q,\Pi]^2)=1$ exactly, so $j^\mu_\Pi=-f^2(\partial^\mu\theta+A^\mu)$, and coupling it to the Maxwell term
reproduces the St\"uckelberg mass $m_A^2=g_A^2f^2$. The mass and the current come from the same trace, so no choice of $f$ gives a massless photon with this source.
\end{theorem}

\begin{proof}
On the ansatz $\Pi=e^{\theta Q}\Pi_0e^{-\theta Q}$ one has $D_\mu\Pi=b_\mu T$ with $b_\mu=\partial_\mu\theta+A_\mu$ and $T=[Q,\Pi]$, and $\Tr(T^2)=1$ is
check g3-7; it is constant along the orbit. By \eqref{eq:jPi}, $j^\mu_\Pi=-f^2b^\mu\Tr([Q,\Pi]T)=-f^2b^\mu\Tr(T^2)=-f^2b^\mu$, and the kinetic term is
$\tfrac{f^2}2\Tr(T^2)\,b_\mu b^\mu=\tfrac{f^2}2b_\mu b^\mu$, the St\"uckelberg Lagrangian of \cref{ex:stueckelberg}. With the Maxwell term, the gauge $\theta=0$
gives the Proca mass $m_A^2=g_A^2f^2$. The current and the mass carry the same factor $f^2\Tr(T^2)$, so a nonzero current requires $f\neq0$ and then gives
a mass.
\end{proof}

\begin{corollary}[The massless criterion]
\label{cor:massless}
A massless vector requires a vacuum $\Pi_0$ with $[Q,\Pi_0]=0$, so that the kinetic term generates no mass, and fluctuations $\pi$ with $[Q,\pi]\neq0$, so that charged modes
still carry the current $-f^2\Tr(Q[\pi,D^\mu\pi])+\cdots$. If $[Q,\Pi(x)]=0$ is imposed everywhere, the current vanishes identically and nothing is sourced.
\end{corollary}

% --- source: tex/s26_photon#3 (The first rung)
\subsection{The first rung}

\begin{theorem}[The criterion can be met]
\label{thm:massless1}
In $\Cl(3,3)$ with $Q=J_N/2$ ($J_N$ of \cref{lem:JN}, which commutes with the four-plane $\langle\widehat W,e_4,e_5,e_6\rangle$ of the symmetric clock) the rank-four orthogonal projectors commuting with $Q$ form $\mathrm{Gr}_\C(2,4)$, of real dimension eight, since $2Q$ is a complex structure.
For a generic member $\Pi_0$,
\[
\dim\{[B,\Pi_0]:B\in\so{3,3}\}=14,\qquad \dim\{\pi\ \text{in that span with}\ [Q,\pi]\neq0\}=8,
\]
so the vacuum is massless and not rigid; the mechanism is the Jacobi identity $[Q,[B,\Pi_0]]=[[Q,B],\Pi_0]$, with $\ad_Q$ of rank eight on $\so{3,3}$. But only two of the eight
charged directions are admissible: by \cref{thm:ghost} a variation $[B,\Pi_0]$ stays among the orthogonal projectors exactly when $B$ is compact, and of the fourteen orbit
directions only five are compact-generated, of the eight charged ones only two.
\end{theorem}

\begin{proposition}[Lorentz invariance makes the first rung rigid]
\label{prop:rigid}
With $\lint$ the interior Lorentz algebra of $\langle\widehat W,e_4,e_5,e_6\rangle$, the commutant of $\{\lint,Q\}$ in $\End(\R^8)$ has dimension four, with symmetric part spanned by $\id$ and $\mathcal V$. So the only rank-four orthogonal projectors commuting
with both the interior Lorentz action and $Q$ are $\tfrac12(\id\pm\mathcal V)$, the chirality projectors of \cref{thm:subdichotomy}. Every bivector commutes with $\mathcal V$, so
these have a spin-orbit tangent space of dimension zero: they are rigid, and by \cref{cor:massless} their current vanishes identically.
\end{proposition}

The paired-mode projector fails the same Lorentz condition, since all six bivectors of $\langle\widehat W,e_4,e_5,e_6\rangle$ fail to commute with it (with the axial clock $e_1$ the pole operator $e_1e_4$ commutes with it, \cref{ch:ch11_curvature}); non-invariance of the observation projector
is a standing feature of the one-constituent construction. What \cref{prop:rigid} adds is that imposing Lorentz invariance together with $[Q,\Pi_0]=0$ forces rigidity
at the first rung. The tower relaxes this: a surplus rotation commuting with the interior generators of every constituent is unique at $n=1$, while at $n=2$ there are
six, and $\ad_Q$ has rank $8$ of $15$ at $n=1$ and $20$ of $66$ at $n=2$ on the bivectors, so the higher rungs offer a genuine choice of which $\uu1$ to gauge.

% --- source: tex/s26_photon#4 (A Lorentz-invariant charged vacuum at the second rung)
\subsection{A Lorentz-invariant charged vacuum at the second rung}
\label{sec:towervac}

\begin{definition}[The second-rung vacuum]
\label{def:towervac}
In $\Cl(6,6)$ with the spectator interior action, whose internal complement has signature $(5,3)$, let $a,b,c,d,e$ be orthonormal positive internal generators: the
surplus pair of the first constituent and the three positive directions of the second. Set
\[
Q:=\tfrac12ab,\qquad H:=abcd,\qquad \Pi_0:=\tfrac12(\id+H).
\]
\end{definition}

\begin{theorem}[Lorentz invariant, uncharged and not rigid]
\label{thm:towervac}
$H^2=\id$, $H^{\!\top}=H$ and $\Tr H=0$, so $\Pi_0$ is an orthogonal projector of rank $32$ on $\R^{64}$. Moreover:
\begin{enumerate}[label=(\roman*),leftmargin=2.2em]
\item $[Q,\Pi_0]=0$, so the kinetic term gives the vacuum no gauge-boson mass;
\item $\Pi_0$ commutes with every interior gamma matrix, hence with the interior Lorentz algebra, and with the internal volume that supplies the complex structure;
\item it is moved by a compact internal bivector: for $B=\tfrac12be$, $\pi=[B,\Pi_0]=-\tfrac12acde$ and $[Q,\pi]=\tfrac12bcde\neq0$, with $\pi$ symmetric and $\Tr(\pi^2)=16$.
\end{enumerate}
\end{theorem}

\begin{proof}
$H$ is a product of four positive generators, so $H^2=(-1)^{4\cdot3/2}\id=\id$ and $H^{\!\top}=H^2H=H$, and it is a traceless blade. $Q=\tfrac12ab$ shares two indices with
$H$ and commutes with it by \eqref{eq:bladecomm}. Each interior gamma shares no index with $H$, and $4\cdot1$ is even, so it commutes. $B=\tfrac12be$ shares one index with
$H$ and anticommutes with it, so $[B,\Pi_0]=\tfrac12[B,H]=BH=\tfrac12\,be\,abcd$; since $b\,e\,a\,b=b^2\,e\,a=e\,a$, this is $\tfrac12\,eacd=-\tfrac12\,acde$. The
blade $acde$ has four positive generators and is symmetric, and $\Tr(\pi^2)=\tfrac14\Tr\id=16$. Finally $Q$ shares one index with $\pi$, so
$[Q,\pi]=2Q\pi=-\tfrac12\,ab\,acde=\tfrac12\,bcde$.
\end{proof}

\begin{theorem}[The compact family and its current]
\label{thm:S4}
Let $I_5:=abcde$ and $Z_i:=I_5v_i$ for $v_i\in\{a,b,c,d,e\}$. The $Z_i$ are symmetric of grade four with $Z_iZ_j+Z_jZ_i=2\delta_{ij}$, so for every unit $n\in\R^5$
\[
\Pi(n):=\tfrac12\bigl(\id+I_5n\bigr)
\]
is a rank-$32$ orthogonal projector commuting with every interior gamma matrix. These form a compact family $S^4=\mathrm{Spin}(5)/\mathrm{Spin}(4)$ with $\Pi_0=\Pi(e)$.
At the vacuum the tangents $T_i=[\tfrac12v_ie,\Pi_0]=\tfrac12I_5v_i$ satisfy $\Tr(T_iT_j)=16\delta_{ij}$ and
\[
[Q,T_a]=-T_b,\qquad [Q,T_b]=T_a,\qquad [Q,T_c]=[Q,T_d]=0 :
\]
one charged complex scalar $\phi=n_a+in_b$ and two neutral real scalars. With $D_\mu\phi=(\partial_\mu-iA_\mu)\phi$ the kinetic term and its current are
\[
\mathcal L_\Pi=8f^2\,D_\mu n\cdot D^\mu n,\qquad j^\mu_\Pi=16f^2\bigl(n_aD^\mu n_b-n_bD^\mu n_a\bigr)=16f^2\Im(\bar\phi D^\mu\phi),
\]
the factor $16$ being the real trace on the $64$-dimensional module. At $n=e$, $\phi=0$, so no $A_\mu A^\mu$ term is generated, while charged fluctuations carry current.
\end{theorem}

\noindent Executable check: a7-3 (exact symbolic or integer arithmetic on $\R^{64}$; \cref{app:methods}); it also checks \cref{thm:towervac}.

\begin{figure}[t]
\centering
\begin{tikzpicture}[>=Stealth,font=\small]
\shade[ball color=tcol!20,opacity=0.9] (0,0) circle (1.9);
\draw[tcol,thick] (0,0) circle (1.9);
\draw[tcol!60,dashed] (-1.9,0) arc (180:360:1.9 and 0.5);
\draw[tcol!60,dotted] (1.9,0) arc (0:180:1.9 and 0.5);
\coordinate (N) at (0,1.9);
\fill[black] (N) circle (2.2pt);
\node[anchor=north,yshift=-4pt] at (N) {$\Pi_0=\Pi(e)$};
\draw[->,scol,very thick] (N) -- ++(1.5,0.35) node[right] {$T_a$};
\draw[->,scol,very thick] (N) -- ++(0.75,0.75) node[above right] {$T_b$};
\draw[scol,->] (0.95,2.35) arc (15:60:0.65);
\node[scol,anchor=west,align=left] at (2.0,3.05) {$Q$ rotates $(T_a,T_b)$:\\ charged scalar $\phi=n_a+in_b$};
\draw[->,gcol,very thick] (N) -- ++(-1.5,0.35) node[left] {$T_c$};
\draw[->,gcol,very thick] (N) -- ++(-0.7,0.8) node[above left] {$T_d$};
\node[gcol,anchor=east,align=right] at (-2.0,3.05) {$[Q,T_c]=[Q,T_d]=0$:\\ two neutral scalars};
\node[align=center] at (0,-2.45) {$S^4=\{\Pi(n)=\tfrac12(\id+I_5n):\ |n|=1\}$, rank $32$, Lorentz invariant};
\end{tikzpicture}
\caption{The compact family of second-rung vacua (\cref{thm:S4}). Every point is a Lorentz-invariant projector; at the vacuum $\Pi_0$, which commutes with the charge $Q$,
the four tangent directions are generated by compact bivectors, two of them rotated into each other by $Q$ and two neutral.}
\label{fig:S4}
\end{figure}

\begin{theorem}[The curved extension]
\label{thm:towercurved}
Reduce the internal connection so that $a,b$ carry the normal $SO(2)$ connection and $c,d,e$ are parallel, $\Omega=\Omega_{\rm Lorentz}+AQ$. Both summands commute with
$\Pi_0$, so
\[
D\Pi_0=0,\qquad [\mathcal R,\Pi_0]=0
\]
for arbitrary spacetime curvature and arbitrary normal curvature, the two conditions that the transport obstruction of \cref{prop:obstruct} separates, here satisfied
together by construction. Throughout the family $\mathcal C_\mu=[\Pi,D_\mu\Pi]=\tfrac14[n,D_\mu n]=\tfrac12\,n\wedge D_\mu n$, a compact internal bivector at every point, so the
adapted connection stays spin-valued away from the vacuum as well.
\end{theorem}

\noindent Executable check: a7-4 (exact symbolic arithmetic; \cref{app:methods}).

For a general projector $[\Pi,D\Pi]$ need not be a bivector: in a normal frame with unit clock $w$ and negative $s,t$, the projector $\Pi=\tfrac12(\id+X)$ with $X=(w+ws)/\sqrt2$
and $B=st$ gives $[\Pi,[B,\Pi]]=\tfrac12(t-st)$, which has a grade-one part. The $S^4$ family avoids this by construction.

\begin{proposition}[The vacuum and the matter]
\label{prop:towermatter}
The stabiliser of $\Pi_0$ in the internal $\so{5,3}$ is $\so4\oplus\so{1,3}$, of dimension $12$ with compact part $\su2\oplus\su2\oplus\su2$, which does not contain $\su3$. The
construction is therefore incompatible with colour at $n=2$; it is compatible from the third rung on:
\begin{center}
\small
\begin{tabular}{@{}cccc@{}}
\toprule
rung & internal algebra & stabiliser of $H$ & compact part\\
\midrule
$n=2$ & $\so{5,3}$, $\dim28$ & $\so4\oplus\so{1,3}$, $\dim12$ & $9$, no $\su3$\\
$n=3$ & $\so{8,6}$, $\dim91$ & $\so4\oplus\so{4,6}$, $\dim51$ & $27$, containing $\so6\supset\su3$\\
$n=4$ & $\so{11,9}$, $\dim190$ & $\so4\oplus\so{7,9}$, $\dim126$ & $63$\\
\bottomrule
\end{tabular}
\end{center}
At $n=3$, the rung of $P_*$, the unbroken compact algebra is $\so4\oplus\so4\oplus\so6$ and contains colour.
\end{proposition}

\noindent A simple $\su3$ has no nonzero homomorphism to $\su2$, so it embeds in no sum of copies of $\su2$. Executable check: a7-5 (exact blade parity;
\cref{app:methods}).

The second-rung vacuum settles that a vacuum which is simultaneously Lorentz invariant, uncharged, not rigid and moved by a compact generator exists, so the first-rung
rigidity of \cref{prop:rigid} is a property of that rung and not a law; its fluctuation spectrum is one charged complex scalar and two neutral ones, its current is an
ordinary charged-scalar current, and its curved extension has $D\Pi_0=0$ with a spin-valued adapted connection across the whole family. It does not settle which point
of $S^4$ is selected, since a spectrum-only potential is constant on the family; whether $j_\Pi$ is the electric current of a fermion, with which charge normalisation;
the values of $f$ and $g_A$; stability; or a gravitational action. It must be done at $n\ge3$ to leave room for colour. And it uses the spectator action throughout:
under the mixed action of \cref{sec:mixed}, $\Pi_0$ is preserved only by the rotations about one axis, its commutators with the six mixed generators span a
five-dimensional space, and the rank-$16$ Dirac block of the coupled model meets the range of $\Pi_0$ only in zero, with $P\Pi_0P$ having eigenvalues $\tfrac14$ and $\tfrac34$
on it. What the coupled model and the second-rung vacuum share is the charge: $Q=\tfrac12J_1$.

The dichotomy left at the end of \cref{sec:transport} is thereby resolved in one direction. If the normal $U(1)$ is to be electromagnetism, the paired-mode projector
must go, and \cref{thm:towervac} shows what can replace it. At one constituent, the paired-mode projector of \cref{ex:pairedmode} fails the photon test.

\begin{remark}[The third rung]
\label{rem:photonn3}
At $n=3$ the shared charge becomes $Q_A=\tfrac12E_2E_3$, the generator of the normal $SO(2)$ connection, and on the chiral block of \cref{sec:pstar} its charge operator
satisfies $-JQ_AP_*=\tfrac12P_*$: it is half the fermion number. Its mixed trace with weak isospin, $\Tr\bigl(P_L(-JQ_A)(-Jt^L_3)^2\bigr)=4$, is nonzero, so with $\su2_L$ gauged
it cannot be gauged on that matter. The trace cancels only on the whole spectator Weyl block, against mirror families that $P_*$ removes. The electric charge
$B_{\rm em}=-\tfrac12E_7E_{13}+\tfrac16H_C$ carries the Standard-Model charges without anomaly. It commutes with four aligned vacua $\Pi_{\pm A},\Pi_{\pm W}$ of the compact
family of vacuum projectors at $n=3$, and it couples to four charged complex scalar fluctuations with the electron's unit of charge. On these counts the normal
$U(1)$ is not electromagnetism at $n=3$, and $B_{\rm em}$ is the charge of the photon candidate (charge identification fixed; the Maxwell term adopted). Vacuum alignment
by a gauged $\su3\oplus\su2\oplus\uu1$ with a positive coefficient selects these four points, and adopted overlaps with the weak-chirality half of the block and with
the family projector select $\Pi_A$. That selection is conditional on the adopted potential and its signs. $\Pi_A$ leaves all of $\su3\oplus\su2\oplus\uu1$ unbroken,
so it is the electroweak-symmetric point. The adopted potential increases along the neutral doublet deformation that would break $\su2\oplus\uu1$ to
$\uu1_{\rm em}$, so a photon distinguished by mass from the $W$ and $Z$, and the mechanism of electroweak breaking, are not obtained. The proofs are in \cref{p1:sec:main,p2:sec:main,p3:sec:main} of this chapter, and the checks are the scripts \texttt{p1}--\texttt{p3} and \texttt{z2} (\cref{app:methods}).
\end{remark}

% --- source: photon_selection/p5_synthesis#0 ((head))

% --- source: photon_selection/p5_synthesis#1 (Synthesis: the photon at the third rung)
\section{Synthesis: the photon at the third rung}
\label{p5:sec:main}

% --- source: photon_selection/p5_synthesis#2 (What is established, adopted and missing)
\subsection{What is established, adopted and missing}

\begin{enumerate}[label=(\arabic*),leftmargin=2.2em]
\item \emph{Charge identification (fixed).} On the chiral block $P_*$, $\hat Q_AP_*=\tfrac12P_*$: the normal-plane charge is half the fermion number and has the
mixed anomaly $\Tr(P_L\hat Q_A(\hat T^3_L)^2)=4$. The electric charge $B_{\rm em}=-\tfrac12E_7E_{13}+\tfrac16H_C$ has the Standard-Model values and no anomaly
(\cref{p3:prop:charges,p3:cor:notem}). The normal $U(1)$ is not electromagnetism on this matter for any vacuum (obstructed). Its anomaly cancels only
together with the mirror families of the spectator Weyl block (\cref{p3:rem:anomaly}). Whether its connection is excluded, cancelled or decoupled is open.
\item \emph{Vacuum family and alignment (fixed, given hypotheses).} The compact family $F_3$ of vacuum projectors carries no selection under $\mathfrak k$-invariant
couplings (\cref{p1:prop:K,p2:prop:allk}). A gauged proper subalgebra with a positive alignment coefficient selects its fixed points
(\cref{p2:prop:mass,p2:prop:zero}). The positivity is the hypothesis (H3), with one-loop support in a cutoff scheme only (\cref{p2:rem:sign}).
\item \emph{Selection of $\Pi_A$ (conditional).} With $\mathfrak h_{\rm SM}$ gauged and $\xi_1>0$ the minima are $\Pi_{\pm A},\Pi_{\pm W}$. At $\xi_1=0$ nothing is
selected (PH-01). Adopted overlaps $-\xi_2\Tr(\Pi P_L)$ and $-\xi_3\Tr(\Pi\Pfam)$ with positive coefficients select $\{\Pi_A,\Pi_{-W}\}$ and then $\Pi_A$. Their
origin and signs are adopted, and the $\xi_3$-term encodes its target (\cref{p3:thm:selection,p3:rem:status}). Local stability needs the conditions of
\cref{p3:prop:hessian}: $\Pi_{-W}$ is not a local minimum when $\xi_1=\xi_2=0<\xi_3$ (PH-02).
\item \emph{Electroweak symmetry (not obtained).} $\Pi_A$ preserves all of $\mathfrak h_{\rm SM}$, so it is the electroweak-symmetric point (PH-03). The neutral
doublet direction $\Pi_\theta$ would give the one-doublet pattern $m_W^2/m_Z^2=(1+\tfrac35g_Y^2/g_2^2)^{-1}$ with $B_{\rm em}$ massless. But the adopted potential
increases along it by $\xi_1 V_{\mathfrak h_{\rm SM}}+(16\xi_2+32\xi_3)(1-\cos\theta)$ (\cref{p3:prop:neutral}). The only other critical set of
$V_{\mathfrak h_{\rm SM}}$ leaves the $Z$ massless (\cref{p2:prop:metastable}). The electroweak mechanism is missing (PH-04).
\item \emph{Photon dynamics (adopted, not derived).} With $B_{\rm em}$ gauged and a Maxwell term adopted, a positive kinetic term exists because the invariant
form is positive on compact generators. The field is massless at the aligned vacua and couples to the block and to four charged complex scalars with
the electron's unit (\cref{p3:thm:photon}). Gauge redundancy, a conserved current, two physical polarisations and quantisation are those of the adopted
Yang--Mills action. A physical photon distinguished from $W^\pm$ and $Z$ requires the missing electroweak stage (PH-06).
\item \emph{Scalars and weak angle (representation content only).} The sixteen tangent directions at an aligned vacuum form four complex doublets. After an
electroweak stage thirteen real scalars would remain, with masses, mixing and Yukawa couplings undetermined (PH-08). The value $\sin^2\theta_W=\tfrac38$
follows from the normalisation $|t^L_3|^2:|t_Y|^2=3:5$ and the assumption $g_Y=g_2$. A prediction would need a matching scale, the spectrum, running and
thresholds (PH-09).
\end{enumerate}

% --- source: photon_selection/p5_synthesis#3 (Two routes to the photon)
\subsection{Two routes to the photon}

Two routes to a Maxwell field have been examined, here and in the global manuscript, and their obstructions are kept.
\begin{itemize}
\item \emph{The normal-plane vector.} Its connection is the $SO(2)$ of the surplus pair. As a Kaluza--Klein vector of a plane of the sign of time it
gravitates with the wrong sign (ghost), it does not couple to spinless matter (\cref{cor:ein6} and \cref{thm:ghostvec} and \cref{prop:constraints} and \cref{thm:wlgauge}), and at $n=3$ its charge is anomalous half fermion
number on the block (obstructed on three counts).
\item \emph{The internal $\uu1_{\rm em}$ and the field-strength reading.} The internal connection along $B_{\rm em}$ with an adopted Maxwell term is healthy and
carries the right charges (constructed, with the charge identification fixed). Separately, the field-strength reading $\R^{3,3}=\Lambda^2\R^{1,3}$ gives the
positive Maxwell energy $|\bB|^2+|\bE|^2$ for a field strength. That reading is kinematic, and connecting it with the internal $\uu1_{\rm em}$ connection, so
that one field is both the gauge field of $B_{\rm em}$ and the field strength of the arena, is not done.
\end{itemize}
A non-preserving projector, a wrong-sign kinetic term or a curvature term of the wrong sign obstructs each construction where the cited theorems
say so; none of these obstructions is removed here.

% --- source: photon_selection/p5_synthesis#4 (Discontinued work)
\subsection{Discontinued work}

The planned track p4, on the gluon if the obstruction persists at $n\ge3$, is cancelled by the author's decision. No gluon identification, gluon-specific
massless-vector, confinement or gluon mass investigation is part of this manuscript. The colour statements that are used, $\su3_C\subset\so6_C$, its trivial
action on $V_+$, and the colour charges of the block, are the established ones of \cref{sec:pstar}.

